\chapter{提升与加性树（Boosting and Additive Trees）}

\textbf{Boosting} 是过去二十年最强大的学习思想之一，最初为分类问题设计，也可扩展至回归。其动机是组合多个「弱」分类器的输出，形成一个强大的「委员会」。从这一角度看，boosting 与 \textbf{Bagging}（§8.8）及其他基于委员会的方法相似；但本质不同——Bagging 是降方差、不扩模型类，而 boosting 是\textbf{通过前向分阶段方式拟合加性模型}，从而扩展模型类。

% ============ §10.1 Boosting 方法 ============
\section{Boosting 方法}

考虑二分类问题，输出编码为 $Y \in \{-1, 1\}$。给定预测变量向量 $X$，分类器 $G(X)$ 输出 $\{-1, 1\}$ 之一。训练样本上的误差率为
\begin{equation}\label{eq:10_err}
    \overline{\mathrm{err}} = \frac{1}{N}\sum_{i=1}^{N} I(y_i \neq G(x_i)),
\end{equation}
未来预测的期望误差率为 $\E_{XY} I(Y \neq G(X))$。

\textbf{弱分类器（weak classifier）}：误差率仅略好于随机猜测的分类器。Boosting 的目标是\textbf{顺序地}把弱分类算法应用到数据的反复修改版本上，产生弱分类器序列 $G_m(x),\, m = 1, 2, \ldots, M$，再通过\textbf{加权多数投票}组合所有预测：
\begin{equation}\label{eq:10_1}
    G(x) = \sgn\Bigl(\sum_{m=1}^{M} \alpha_m G_m(x)\Bigr),
\end{equation}
其中 $\alpha_1, \ldots, \alpha_M$ 由 boosting 算法计算，衡量各 $G_m(x)$ 的贡献；其作用是给序列中更准确的分类器更高影响力。整个流程可示意如下：

\begin{center}
\begin{tikzpicture}[>=stealth, font=\small,
    dbox/.style={draw, rounded corners=2pt, inner sep=4pt, align=center, fill=blue!8},
    cbox/.style={draw, rounded corners=2pt, inner sep=4pt, align=center, fill=orange!12},
    fbox/.style={draw, rounded corners=2pt, inner sep=5pt, align=center, fill=green!10, thick}]
  % 数据列：训练样本 → 反复修改的加权样本
  \node[dbox] (tr) at (0,0)    {训练样本};
  \node[dbox] (d1) at (0,-1.7) {加权样本$_1$};
  \node[dbox] (d2) at (0,-3.4) {加权样本$_2$};
  \node[dbox] (d3) at (0,-5.1) {加权样本$_3$};
  \node at (0,-6.2) {$\vdots$};
  \node[dbox] (dM) at (0,-7.3) {加权样本$_M$};
  % 分类器列
  \node[cbox] (g1) at (3.7,-1.7) {$G_1(x)$};
  \node[cbox] (g2) at (3.7,-3.4) {$G_2(x)$};
  \node[cbox] (g3) at (3.7,-5.1) {$G_3(x)$};
  \node at (3.7,-6.2) {$\vdots$};
  \node[cbox] (gM) at (3.7,-7.3) {$G_M(x)$};
  % 最终分类器
  \node[fbox] (F) at (7.8,-4.4) {最终分类器\\[2pt] $G(x) = \sgn\Bigl(\sum_{m=1}^{M} \alpha_m G_m(x)\Bigr)$};
  % 数据 → 加权样本（更新权重）
  \draw[->, thick] (tr) -- (d1) node[midway, right] {初始 $w_i = 1/N$};
  \draw[->] (d1) -- (d2) node[midway, right] {更新权重 $w_i$};
  \draw[->] (d2) -- (d3);
  \draw[->, dashed] (d3) -- (dM);
  % 加权样本 → 弱分类器
  \draw[->] (d1) -- (g1);
  \draw[->] (d2) -- (g2);
  \draw[->] (d3) -- (g3);
  \draw[->] (dM) -- (gM);
  % 弱分类器 → 最终分类器（加权多数投票）
  \draw[->] (g1) -- (F) node[midway, above, sloped] {$\alpha_1$};
  \draw[->] (g2) -- (F) node[midway, above, sloped] {$\alpha_2$};
  \draw[->] (g3) -- (F) node[midway, above, sloped] {$\alpha_3$};
  \draw[->] (gM) -- (F) node[midway, below, sloped] {$\alpha_M$};
\end{tikzpicture}
\\[2pt]
\small\textbf{图：}AdaBoost 示意：对训练数据的加权版本顺序训练弱分类器序列 $G_1(x), \ldots, G_M(x)$，每步把被 $G_m$ 误分类的观测权重放大，再以权重 $\alpha_m$ 进行加权多数投票，组合成最终分类器 $G(x) = \sgn\bigl(\sum_{m=1}^{M} \alpha_m G_m(x)\bigr)$。
\end{center}

\paragraph{sgn 怎么理解？}

$\sgn(\cdot)$ 是\textbf{符号函数}：
\[
    \sgn(z) = \begin{cases} 1, & z > 0,\\ -1, & z < 0,\\ 0, & z = 0. \end{cases}
\]
在式 \eqref{eq:10_1} 中，每个弱分类器 $G_m(x) \in \{-1, 1\}$ 相当于对类标签\textbf{投一票}（$+1$ 投给类 $1$、$-1$ 投给类 $-1$），$\alpha_m > 0$ 是这一票的权重；于是 $\sum_{m=1}^{M} \alpha_m G_m(x)$ 就是\textbf{加权投票的净得分}——所有投 $+1$ 的权重之和减去所有投 $-1$ 的权重之和。\textbf{sgn 的作用}是把这个实值净得分翻译成最终类标签：净得分为正判为类 $1$、为负判为类 $-1$，即
\[
    G(x) = +1 \iff \sum_{m:\,G_m(x)=+1} \alpha_m > \sum_{m:\,G_m(x)=-1} \alpha_m.
\]
所以 $G(x) = \sgn\bigl(\sum_{m=1}^{M} \alpha_m G_m(x)\bigr)$ 的直观就是\textbf{加权多数投票}（weighted majority vote）。注意 sgn 只保留正负号、丢掉幅度信息——训练阶段 AdaBoost 实际拟合的是\textbf{连续值} $f(x) = \sum_{m=1}^{M} \beta_m G_m(x)$（§10.4），指数损失 $\exp(-y f(x))$ 依赖连续的 $f$ 而非其符号；sgn 只是最后一步的决策读出。

数据修改方式：在每一步 boosting 中，对每个训练观测 $(x_i, y_i)$ 施加权重 $w_1, w_2, \ldots, w_N$。初始设 $w_i = 1/N$（第一步即在原始数据上照常训练）。对后续迭代 $m = 2, 3, \ldots, M$，被前一步分类器 $G_{m-1}(x)$ 误分类的观测权重\textbf{增加}，正确分类的观测权重\textbf{减少}。于是随着迭代进行，难以正确分类的观测获得越来越大的影响力，迫使每个后续分类器专注于先前分类器遗漏的样本。

\textbf{Algorithm 10.1（AdaBoost.M1，离散 AdaBoost）：}
\begin{enumerate}
    \item 初始化观测权重 $w_i = 1/N,\; i = 1, 2, \ldots, N$。
    \item 对 $m = 1$ 到 $M$：
    \begin{enumerate}
        \item 用权重 $w_i$ 在训练数据上拟合分类器 $G_m(x)$。
        \item 计算加权误差率
        \begin{equation}\label{eq:10_errm}
            \mathrm{err}_m = \frac{\sum_{i=1}^{N} w_i I(y_i \neq G_m(x_i))}{\sum_{i=1}^{N} w_i}.
        \end{equation}
        \item 计算 $\alpha_m = \log\bigl((1 - \mathrm{err}_m)/\mathrm{err}_m\bigr)$。
        \item 更新权重 $w_i \leftarrow w_i \cdot \exp\bigl[\alpha_m \cdot I(y_i \neq G_m(x_i))\bigr],\; i = 1, 2, \ldots, N$。
    \end{enumerate}
    \item 输出 $G(x) = \sgn\bigl[\sum_{m=1}^{M} \alpha_m G_m(x)\bigr]$。
\end{enumerate}

第 2(c) 行计算的 $\alpha_m$ 是最终分类器中给 $G_m$ 的权重：误分类观测的权重被放大 $\exp(\alpha_m)$ 倍，增强其对下一步分类器 $G_{m+1}$ 的影响力。

\textbf{离散 AdaBoost}：因基分类器 $G_m(x)$ 返回离散类标签而得名（Friedman et al., 2000）。若基分类器改为返回实值预测（例如映射到 $[-1,1]$ 区间的概率），可相应修改得到「Real AdaBoost」。

\begin{example}[模拟数据：stump 的威力]
特征 $X_1, \ldots, X_{10}$ 为标准独立高斯，确定性目标
\begin{equation}\label{eq:10_2}
    Y = \begin{cases} 1, & \text{若 } \sum_{j=1}^{10} X_j^2 > \chi^2_{10}(0.5),\\ -1, & \text{否则}, \end{cases}
\end{equation}
其中 $\chi^2_{10}(0.5) = 9.34$ 是 10 自由度卡方随机变量的中位数（10 个标准高斯的平方和）。2000 个训练样本（每类约 1000 个），10000 个测试观测。弱分类器是「stump」——仅含两个终端节点的分类树。
\begin{itemize}
    \item 单独应用 stump：测试误差 $45.8\%$（随机猜测为 $50\%$），几乎无效。
    \item 400 次 boosting 迭代后：误差降至 $5.8\%$，约为原误差率的 $\frac{1}{4}$。
    \item 单个 244 节点的大分类树：误差率 $24.7\%$——boosting 完胜。
\end{itemize}
Breiman（NIPS Workshop, 1996）称「带树的 AdaBoost 是世界上最好的现成分类器」（见 Breiman, 1998），对数据挖掘应用尤其如此（§10.7 详述）。
\end{example}

\subsection{本章概览}

本章的要点可概述如下：
\begin{itemize}
    \item AdaBoost 在基学习器上拟合\textbf{加性模型}，优化一种新颖的\textbf{指数损失}，该损失与（负）二项对数似然非常相似（§10.2--§10.4）。
    \item 指数损失的\textbf{总体极小化}被证明是类概率的对数优势（§10.5）。
    \item 给出比平方误差或指数损失更稳健的回归与分类损失函数（§10.6）。
    \item 论证决策树是数据挖掘应用中 boosting 的理想基学习器（§10.7、§10.9）。
    \item 发展一类\textbf{梯度提升模型（GBM）}，用于任意损失下对树进行 boosting（§10.10）。
    \item 强调「慢学习」的重要性，通过每个新进入模型的项的\textbf{收缩（shrinkage）}（§10.12）及随机化（§10.12.2）实现。
    \item 描述拟合模型的解释工具（§10.13）。
\end{itemize}

% ============ §10.2 Boosting 拟合加性模型 ============
\section{Boosting 拟合加性模型}

Boosting 的成功并不神秘，关键在式 \eqref{eq:10_1}：Boosting 是在一组初等「基」函数上\textbf{拟合加性展开}的方法。此处基函数即单个分类器 $G_m(x) \in \{-1, 1\}$。更一般地，基函数展开形如
\begin{equation}\label{eq:10_3}
    f(x) = \sum_{m=1}^{M} \beta_m b(x; \gamma_m),
\end{equation}
其中 $\beta_m$ 是展开系数，$b(x; \gamma) \in \R$ 通常是多元自变量 $x$ 的简单函数，由参数集 $\gamma$ 刻画（基展开详见第 5 章）。

这类加性展开是书中许多学习技术的核心：
\begin{itemize}
    \item \textbf{单隐层神经网络}（第 11 章）：$b(x; \gamma) = \sigma(\gamma_0 + \gamma_1^T x)$，其中 $\sigma(t) = 1/(1 + e^{-t})$ 是 sigmoid 函数，$\gamma$ 参数化输入变量的线性组合。
    \item \textbf{小波}（§5.9.1）：$\gamma$ 参数化「母小波」的位置与尺度平移。
    \item \textbf{MARS}（§9.4）：截断幂样条基，$\gamma$ 参数化结点变量与结点值。
    \item \textbf{树}：$\gamma$ 参数化内部结点的分裂变量与分裂点、以及终端结点的预测。
\end{itemize}

这类模型通常通过最小化训练数据上的平均损失（如平方误差或基于似然的损失）来拟合：
\begin{equation}\label{eq:10_4}
    \min_{\{\beta_m, \gamma_m\}_1^M} \frac{1}{N}\sum_{i=1}^{N} L\Bigl(y_i, \sum_{m=1}^{M} \beta_m b(x_i; \gamma_m)\Bigr).
\end{equation}
对许多损失 $L(y, f(x))$ 与基函数 $b(x; \gamma)$，这需要计算密集的数值优化技术。但若可以\textbf{快速求解仅拟合单个基函数的子问题}，则存在简单替代方案：
\begin{equation}\label{eq:10_5}
    \min_{\beta, \gamma} \sum_{i=1}^{N} L\bigl(y_i, \beta b(x_i; \gamma)\bigr).
\end{equation}

\paragraph{Boosting 拟合的「加性模型」是基于 GAM 吗？}

\textbf{不是}。两者虽然都叫「加性」，但「加性」的含义不同：
\begin{itemize}
    \item \textbf{GAM（第 9 章）}：加性是\textbf{按特征}加性——$f(x) = \alpha + \sum_{j=1}^{p} f_j(X_j)$，每个预测变量 $X_j$ 贡献一个光滑函数 $f_j$，各项是\textbf{单个变量}的函数（主效应），交互需显式加入（如 $f_{jk}(X_j, X_k)$）。
    \item \textbf{Boosting}：加性是\textbf{按基学习器}加性——$f(x) = \sum_{m=1}^{M} \beta_m b(x; \gamma_m)$，每个基函数 $b(x; \gamma_m)$ 都是\textbf{整个输入向量 $x$} 的函数（如完整树、完整 stump），可同时依赖任意特征组合，因而\textbf{天然包含交互}。
\end{itemize}
更准确的定位：两者都是「\textbf{基函数加性展开}」这一更一般框架（§10.2 开头视角，与第 5 章基展开同构）的特例，只是\textbf{基函数的选择与拟合方式不同}：
\begin{itemize}
    \item GAM 的基 = 预定义/数据驱动的\textbf{单变量光滑函数}，用 \textbf{backfitting}（交替更新，§9.1.1）\textbf{同时}拟合所有项；
    \item Boosting 的基 = \textbf{自适应地从数据中学到的弱学习器}（每步贪心地选最好的基加入），用\textbf{前向分阶段}（Algorithm 10.2）\textbf{顺序}拟合——每次只加一项、不回改已添加的项。
\end{itemize}
所以不能把 boosting 理解成「基于 GAM」；更准确的说法是：boosting 用前向分阶段方式拟合「弱学习器基函数」上的加性展开，而 GAM 用 backfitting 拟合「特征光滑函数」上的加性展开。二者仅共享「把若干项加起来」的展开形式；分道扬镳之处在于\textbf{项长什么样}（单变量光滑函数 vs 全特征弱学习器）以及\textbf{怎么拟合}（同时交替更新 vs 顺序贪心添加）。这也解释了为什么 boosting 能自动捕捉交互而纯 GAM 不能——每项（如树）本身就可以是特征的非线性 + 交互函数。

% ============ §10.3 前向分阶段加性建模 ============
\section{前向分阶段加性建模}

\textbf{前向分阶段建模（forward stagewise modeling）}通过\textbf{顺序添加}新基函数来逼近式 \eqref{eq:10_4} 的解，\textbf{不调整已添加项的参数与系数}（Algorithm 10.2）。每步迭代 $m$，求解要加入当前展开 $f_{m-1}(x)$ 的最优基函数 $b(x; \gamma_m)$ 及其系数 $\beta_m$，得到 $f_m(x)$ 后重复；之前添加的项不再修改。

\textbf{Algorithm 10.2（前向分阶段加性建模，Forward Stagewise Additive Modeling）：}
\begin{enumerate}
    \item 初始化 $f_0(x) = 0$。
    \item 对 $m = 1$ 到 $M$：
    \begin{enumerate}
        \item 计算 $(\beta_m, \gamma_m) = \argmin_{\beta, \gamma} \sum_{i=1}^{N} L\bigl(y_i, f_{m-1}(x_i) + \beta b(x_i; \gamma)\bigr)$。
        \item 令 $f_m(x) = f_{m-1}(x) + \beta_m b(x; \gamma_m)$。
    \end{enumerate}
\end{enumerate}

对\textbf{平方误差损失}
\begin{equation}\label{eq:10_6}
    L(y, f(x)) = (y - f(x))^2,
\end{equation}
有
\begin{equation}\label{eq:10_7}
    L\bigl(y_i, f_{m-1}(x_i) + \beta b(x_i; \gamma)\bigr) = (y_i - f_{m-1}(x_i) - \beta b(x_i; \gamma))^2 = (r_{im} - \beta b(x_i; \gamma))^2,
\end{equation}
其中 $r_{im} = y_i - f_{m-1}(x_i)$ 恰是当前模型对第 $i$ 个观测的\textbf{残差}。因此，平方误差损失下，每步加入的项 $\beta_m b(x; \gamma_m)$ 就是\textbf{对当前残差拟合最好的项}。这一思想是 §10.10.2「最小二乘回归提升」的基础。但如后文所示，平方误差损失一般不是分类的好选择，因此需要考虑其他损失准则。

% ============ §10.4 指数损失与 AdaBoost ============
\section{指数损失与 AdaBoost}

下面证明 AdaBoost.M1（Algorithm 10.1）\textbf{等价于}使用损失
\begin{equation}\label{eq:10_8}
    L(y, f(x)) = \exp(-y f(x))
\end{equation}
的前向分阶段加性建模（Algorithm 10.2）。对 AdaBoost，基函数即单个分类器 $G_m(x) \in \{-1, 1\}$。使用指数损失，每步需求解分类器 $G_m$ 与相应系数 $\beta_m$：
\begin{equation}\label{eq:10_9}
    (\beta_m, G_m) = \argmin_{\beta, G} \sum_{i=1}^{N} w_i^{(m)} \exp\bigl(-\beta y_i G(x_i)\bigr), \qquad w_i^{(m)} = \exp\bigl(-y_i f_{m-1}(x_i)\bigr).
\end{equation}
由于每个 $w_i^{(m)}$ 既不依赖 $\beta$ 也不依赖 $G(x)$，可视为施加到每个观测上的\textbf{权重}；该权重依赖 $f_{m-1}(x_i)$，故随迭代 $m$ 变化。

式 \eqref{eq:10_9} 的解分两步获得。

\textbf{第一步}：对任意 $\beta > 0$，关于 $G_m$ 的最优解为
\begin{equation}\label{eq:10_10}
    G_m = \argmin_{G} \sum_{i=1}^{N} w_i^{(m)} I(y_i \neq G(x_i)),
\end{equation}
即\textbf{最小化加权误差率}的分类器。这是因为式 \eqref{eq:10_9} 的准则可写成
\begin{equation}\label{eq:10_11}
    (e^{\beta} - e^{-\beta}) \cdot \sum_{i=1}^{N} w_i^{(m)} I(y_i \neq G(x_i)) + e^{-\beta} \cdot \sum_{i=1}^{N} w_i^{(m)}.
\end{equation}

\textbf{第二步}：把第一步得到的 $G_m$ \textbf{固定}，代回式 \eqref{eq:10_9}，把 $\beta$ 解出来。此时 $G(x_i)$ 已知，只需关于 $\beta$ 极小化。先把求和按「分类是否正确」拆成两部分——注意对 $y_i \in \{\pm 1\}$、$G_m(x_i) \in \{\pm 1\}$，恒有 $y_i G_m(x_i) = 1$（正确分类）或 $-1$（误分类）：
\begin{itemize}
    \item 若 $y_i = G_m(x_i)$（正确），则 $y_i G_m(x_i) = 1$，该项贡献 $w_i^{(m)} e^{-\beta}$；
    \item 若 $y_i \neq G_m(x_i)$（误分类），则 $y_i G_m(x_i) = -1$，该项贡献 $w_i^{(m)} e^{\beta}$。
\end{itemize}
于是准则 \eqref{eq:10_9} 化为
\[
    Q(\beta) = e^{-\beta}\underbrace{\sum_{y_i = G_m(x_i)} w_i^{(m)}}_{W^+} + e^{\beta}\underbrace{\sum_{y_i \neq G_m(x_i)} w_i^{(m)}}_{W^-},
\]
其中 $W^+$、$W^-$ 分别为正确分类、误分类观测的\textbf{权重和}，且 $W^+ + W^- = \sum_{i=1}^{N} w_i^{(m)}$。注意 $W^+$、$W^-$ 均\textbf{与 $\beta$ 无关}（$w_i^{(m)}$ 只依赖 $f_{m-1}$），故可按普通常数对 $\beta$ 求导。

对 $\beta$ 求导并令其为零：
\[
    \frac{dQ}{d\beta} = -W^+ e^{-\beta} + W^- e^{\beta} \overset{!}{=} 0
    \;\Longrightarrow\;
    W^- e^{\beta} = W^+ e^{-\beta}
    \;\Longrightarrow\;
    e^{2\beta} = \frac{W^+}{W^-},
\]
两边取对数得
\[
    \beta = \frac{1}{2}\log\frac{W^+}{W^-}.
\]
再把加权误差率 \eqref{eq:10_13} 用进来：$\mathrm{err}_m = \frac{W^-}{W^+ + W^-}$，故 $W^+ = (1 - \mathrm{err}_m)(W^+ + W^-)$、$W^- = \mathrm{err}_m (W^+ + W^-)$，代入即得
\begin{equation}\label{eq:10_12}
    \beta_m = \frac{1}{2}\log\frac{1 - \mathrm{err}_m}{\mathrm{err}_m}.
\end{equation}
其中 $\mathrm{err}_m$ 为最小化后的加权误差率
\begin{equation}\label{eq:10_13}
    \mathrm{err}_m = \frac{\sum_{i=1}^{N} w_i^{(m)} I(y_i \neq G_m(x_i))}{\sum_{i=1}^{N} w_i^{(m)}}.
\end{equation}

\textbf{这是极小点吗？} 二阶导数 $Q''(\beta) = W^+ e^{-\beta} + W^- e^{\beta} > 0$ 恒正，故 $Q(\beta)$ 是 $\beta$ 的\textbf{凸函数}，上述驻点是\textbf{唯一全局极小点}。此外，$\beta_m > 0$ 当且仅当 $\mathrm{err}_m < \frac{1}{2}$——这正是「\textbf{弱分类器}（略优于随机猜测）」假设所保证的：若基分类器误差率超过 $\frac{1}{2}$，则 $\beta_m \le 0$，权重更新方向反转，boosting 失效。这也解释了为何 $w_i^{(m)}$ 与 $\beta$ 无关是推导的关键——它把式 \eqref{eq:10_9} 的联合极小化化成了「先找使加权误差率最小的 $G_m$，再闭式求 $\beta_m$」的两步。\textbf{直觉}：$\beta_m$ 是「对 $G_m$ 有多信任」的度量——$G_m$ 越准（$\mathrm{err}_m \to 0$），$\beta_m$ 越大，最终投票 \eqref{eq:10_1} 中其权重越高；纯随机（$\mathrm{err}_m = \frac{1}{2}$）时 $\beta_m = 0$，该分类器不参与投票。

随后更新逼近 $f_m(x) = f_{m-1}(x) + \beta_m G_m(x)$，使下一步权重变为
\begin{equation}\label{eq:10_14}
    w_i^{(m+1)} = w_i^{(m)} \cdot e^{-\beta_m y_i G_m(x_i)}.
\end{equation}
利用恒等式 $-y_i G_m(x_i) = 2 \cdot I(y_i \neq G_m(x_i)) - 1$，式 \eqref{eq:10_14} 化为
\begin{equation}\label{eq:10_15}
    w_i^{(m+1)} = w_i^{(m)} \cdot e^{\alpha_m I(y_i \neq G_m(x_i))} \cdot e^{-\beta_m},
\end{equation}
其中 $\alpha_m = 2\beta_m$ 正是 AdaBoost.M1 第 2(c) 行定义的量。因子 $e^{-\beta_m}$ 对所有权重乘以同一常数，无影响。因此式 \eqref{eq:10_15} 等价于 Algorithm 10.1 的第 2(d) 行。

\begin{remark}
Algorithm 10.1 的第 2(a) 行可视为对式 \eqref{eq:10_11}（从而 \eqref{eq:10_10}）极小化的近似求解方法。故结论：\textbf{AdaBoost.M1 通过前向分阶段加性建模最小化指数损失准则 \eqref{eq:10_8}}。
\end{remark}

\begin{remark}[AdaBoost 不优化训练集误分类误差]
对式 \eqref{eq:10_2} 的模拟数据，训练集误分类误差约在 250 次迭代后降到零（并保持为零），但平均指数损失 $\frac{1}{N}\sum_i \exp(-y_i f(x_i))$ 仍在持续下降；同时测试集误分类误差在 250 次迭代后仍继续改善。这说明 AdaBoost 优化的并非训练集误分类误差——指数损失对估计类概率的变化更敏感。
\end{remark}

% ============ §10.5 为什么用指数损失 ============
\section{为什么用指数损失}

AdaBoost.M1 最初并非从上一节的角度提出，其与基于指数损失的前向分阶段加性建模的等价性是提出五年后才被发现的。研究指数损失准则的性质，可以深入理解该过程并发现改进方向。

指数损失在加性建模语境中的\textbf{主要吸引力是计算上的}：它导致简单、模块化的重新加权 AdaBoost 算法。但值得追问其统计性质：它估计什么？估计得好不好？第一个问题由求其\textbf{总体极小化}回答。可以证明（Friedman et al., 2000）：
\begin{equation}\label{eq:10_16}
    f^*(x) = \argmin_{f(x)} \E_{Y \mid x}\bigl(e^{-Y f(x)}\bigr) = \frac{1}{2}\log\frac{\Pr(Y = 1 \mid x)}{\Pr(Y = -1 \mid x)},
\end{equation}
即指数损失的总体极小化正是\textbf{类概率的对数优势（log-odds）的一半}。等价地
\begin{equation}\label{eq:10_17}
    \Pr(Y = 1 \mid x) = \frac{1}{1 + e^{-2 f^*(x)}}.
\end{equation}
因此 AdaBoost 估计的是（对数优势意义下的）类概率，而非简单的分类边界。

\begin{remark}[边际与损失函数]
记\textbf{边际} $y f(x)$。分类算法的目标是尽可能产生正边际。任何分类损失都应对负边际惩罚更重、对正边际奖励更轻。误分类损失 $I(yf(x) < 0)$ 对负边际罚 1、对正边际不罚；指数损失与二项偏差（deviance）都是它的单调连续逼近，且对越来越负的边际惩罚比奖励更重——区别只在程度上：二项偏差对大负边际的惩罚\textbf{线性}增长，而指数损失对这类观测的影响力\textbf{指数}增长，因此在训练过程中把更多影响力集中在负边际大的观测上。
\end{remark}

% ============ §10.6 损失函数与稳健性 ============
\section{损失函数与稳健性}

本节系统比较二分类的各种损失函数，并给出更稳健的替代准则。记响应 $y \in \{\pm 1\}$、预测为 $f(x)$、类预测为 $\sgn(f(x))$，\textbf{边际}为 $y f(x)$。分类算法的目标是尽可能产生正边际；任何分类损失都应\textbf{对负边际惩罚更重}（负边际观测已被误分类）、对正边际奖励更轻。

\textbf{误分类损失} $L(y, f(x)) = I(yf(x) < 0)$：对负边际罚 1、对正边际不罚。指数损失与二项偏差（deviance）都可看作它的\textbf{单调连续逼近}，区别只在程度：
\begin{itemize}
    \item \textbf{指数损失} $L(y, f) = e^{-yf}$（式 \eqref{eq:10_8}）：对越来越负的边际惩罚呈\textbf{指数}增长，把更多影响力集中在负边际大的观测上——在含噪（贝叶斯误差不为零）或训练标签被错误标注的情形下，性能会显著下降。
    \item \textbf{二项偏差} $L(y, f) = \log\bigl(1 + e^{-2yf}\bigr)$（负二项对数似然）：对大负边际的惩罚近似\textbf{线性}增长，影响力更均匀地分布于全部数据，因而\textbf{稳健得多}。
\end{itemize}

\paragraph{二项偏差 $\log(1+e^{-2yf})$ 真的是「负二项对数似然」吗？}

\textbf{是}。它正是把类概率写成 logistic 形式后的负对数似然。设 $p = \Pr(Y=1 \mid x)$，取 logistic 形式
\[
    p = \frac{1}{1+e^{-2f}}, \qquad \Pr(Y=-1 \mid x) = 1-p = \frac{1}{1+e^{2f}},
\]
（这里用 $2f$ 是因为 $f$ 定义为 log-odds 的\textbf{一半}，见式 \eqref{eq:10_16}--\eqref{eq:10_17}）。对单个观测的负对数似然为
\[
    -\log \Pr(Y=y \mid x) = \begin{cases} \log(1+e^{-2f}), & y = 1,\\ \log(1+e^{2f}), & y = -1, \end{cases}
\]
由于 $y \in \{\pm 1\}$ 时 $-2yf$ 分别等于 $-2f$ 与 $+2f$，两种情形可统一写成
\[
    L(y, f) = -\log \Pr(Y = y \mid x) = \log\bigl(1 + e^{-2yf}\bigr),
\]
即\textbf{负二项对数似然}（二项偏差）。它也叫「logistic 损失」。它与指数损失 $e^{-yf}$ 的关系：对 $\log(1+e^{-2yf})$ 做大负边际近似得 $2\cdot(-yf) = 2|yf|$（线性），而指数损失近似得 $e^{|yf|}$（指数）——这正是 §10.6 所说「二项偏差对大负边际惩罚线性、指数损失指数增长」的由来。

三个损失函数随边际 $yf(x)$ 的变化如图：
\begin{center}
\begin{tikzpicture}[>=stealth, font=\small]
    % 坐标轴与刻度（先画，避免被 clip 裁剪）
    \draw[->, thick] (-2.6,0) -- (2.7,0) node[below right] {$yf(x)$（边际）};
    \draw[->, thick] (0,-0.12) -- (0,3.45) node[above left] {损失};
    \draw (-2,0.05) -- (-2,-0.05) node[below] {$-2$};
    \draw (-1,0.05) -- (-1,-0.05) node[below] {$-1$};
    \draw (1,0.05) -- (1,-0.05) node[below] {$1$};
    \draw (2,0.05) -- (2,-0.05) node[below] {$2$};
    \draw (0.05,1) -- (-0.05,1) node[left] {$1$};
    \draw (0.05,2) -- (-0.05,2) node[left] {$2$};
    \draw (0.05,3) -- (-0.05,3) node[left] {$3$};
    % 曲线（clip 在坐标轴范围内，允许函数超出顶部被截断）
    \begin{scope}
        \clip (-2.6,-0.2) rectangle (2.6,3.4);
        % 误分类损失：阶梯（yf<0 取 1，yf>=0 取 0）
        \draw[very thick, black] (-2.5,1) -- (0,1);
        \draw[very thick, black] (0,0) -- (2.5,0);
        % 指数损失 exp(-yf)，0 处 =1，无需缩放
        \draw[red, thick, domain=-2.5:2.5, samples=200] plot (\x, {exp(-\x)});
        % 二项偏差 log(1+e^{-2yf})/log 2，缩放使其过 (0,1)
        \draw[blue, thick, domain=-2.5:2.5, samples=200] plot (\x, {ln(1+exp(-2*\x))/ln(2)});
    \end{scope}
    % 误分类损失在 yf=0 处：I(yf<0)=0，故 (0,0) 实心、(0,1) 空心
    \filldraw[black] (0,0) circle (1.6pt);
    \filldraw[white] (0,1) circle (1.6pt);
    \draw (0,1) circle (1.6pt);
    % 图例（右上角空白区）
    \draw[black, very thick] (1.05,2.95) -- (1.45,2.95) node[right, black] {误分类 $I(yf<0)$};
    \draw[red, thick] (1.05,2.6) -- (1.45,2.6) node[right, red] {指数 $e^{-yf}$};
    \draw[blue, thick] (1.05,2.25) -- (1.45,2.25) node[right, blue] {二项偏差 $\log(1+e^{-2yf})$};
\end{tikzpicture}
\\[2pt]
\small\textbf{图：}二分类的三种损失函数（缩放后均过 $(0,1)$ 点，对应原书图 10.4）。误分类损失是「对负边际罚 1、正边际不罚」的阶梯；指数损失与二项偏差都是它的\textbf{单调连续逼近}，且都随边际变负而增大——指数损失上升\textbf{指数级}（对误分类、尤其是标签噪声的观测施加极大影响力），二项偏差上升近似\textbf{线性}（影响力更均匀、更稳健）。
\end{center}

\subsection{平方误差损失不适合分类}

平方误差损失 $L(y, f) = (y - f)^2$ 的总体风险极小化者为
\begin{equation}\label{eq:10_19}
    f^*(x) = \argmin_{f(x)} \E_{Y \mid x}(Y - f(x))^2 = \E(Y \mid x) = 2\cdot \Pr(Y = 1 \mid x) - 1,
\end{equation}
分类规则仍取 $G(x) = \sgn[f(x)]$。

\paragraph{$\E(Y \mid x) = 2\Pr(Y = 1 \mid x) - 1$ 这个等号怎么成立的？}

这完全来自二分类的\textbf{编码约定} $Y \in \{-1, 1\}$。设 $p = \Pr(Y = 1 \mid x)$，则 $\Pr(Y = -1 \mid x) = 1 - p$。按（离散）期望的定义把 $Y$ 的两个取值 $1$ 与 $-1$ 及其条件概率代入：
\[
    \E(Y \mid x) = \sum_y y \cdot \Pr(Y = y \mid x)
    = 1 \cdot p + (-1) \cdot (1 - p)
    = p - 1 + p
    = 2p - 1
    = 2\Pr(Y = 1 \mid x) - 1.
\]
\textbf{直觉}：$Y$ 的两个取值 $+1$、$-1$ 关于 $0$ 对称，因此条件期望就是「两类概率之差的加权平衡点」——当 $\Pr(Y=1\mid x) = \frac{1}{2}$（两类等可能）时 $\E(Y\mid x) = 0$（在 $0$ 处平衡）；当类 $1$ 概率趋向 $1$ 时，期望趋向 $+1$；趋向 $0$ 时趋向 $-1$。它把区间 $[0,1]$ 上的概率线性映射到 $[-1,1]$ 上的期望，故 $\sgn[\E(Y\mid x)] = \sgn[2p-1]$ 恰好实现贝叶斯分类规则（$p > \frac{1}{2}$ 判类 $1$）。

但平方误差\textbf{不是误分类误差的好替代}：它不是边际 $yf(x)$ 的单调递减函数——当 $y_i f(x_i) > 1$ 时它反而\textbf{二次增大}，把越来越多误差加在「越来越确定被正确分类」的观测上，从而\textbf{相对削弱}了真正误分类（$y_i f(x_i) < 0$）观测的影响力。若目标是类别判定，应采用单调递减的替代损失。第 12 章图 12.4 给出二次损失的修改——「Huber 化的平方 hinge 损失」（Rosset et al., 2004b），兼具二项偏差、二次损失与 SVM hinge 损失的良好性质：与二次损失 \eqref{eq:10_19} 有相同的总体极小化、$yf(x) > 1$ 时取零、$yf(x) < -1$ 时变为线性。

\subsection{K 类分类的损失}

响应 $Y$ 取值于无序集合 $\mathcal{G} = \{\mathcal{G}_1, \ldots, \mathcal{G}_K\}$（§2.4、§4.4）。只需类条件概率 $p_k(x) = \Pr(Y = \mathcal{G}_k \mid x)$ 即可构造贝叶斯分类器
\begin{equation}\label{eq:10_20}
    G(x) = \mathcal{G}_k, \quad \text{其中 } k = \argmax_{\ell} p_\ell(x).
\end{equation}
数据挖掘应用中常更关心类概率本身。如 §4.4，logistic 模型自然推广到 $K$ 类：
\begin{equation}\label{eq:10_21}
    p_k(x) = \frac{e^{f_k(x)}}{\sum_{\ell=1}^{K} e^{f_\ell(x)}},
\end{equation}
保证 $0 \le p_k(x) \le 1$ 且求和为 1。这里每类一个函数 $f_k$；因给所有 $f_k$ 加上同一 $h(x)$ 不改变模型，存在冗余。传统上令其中一个为零（如 $f_K(x) = 0$，见 (4.17)）；本书偏好保留对称性，施加约束 $\sum_{k=1}^{K} f_k(x) = 0$。二项偏差自然推广为 \textbf{K 类多项式偏差}：
\begin{equation}\label{eq:10_22}
    L(y, p(x)) = -\sum_{k=1}^{K} I(y = \mathcal{G}_k)\log p_k(x) = -\sum_{k=1}^{K} I(y = \mathcal{G}_k) f_k(x) + \log\Bigl(\sum_{\ell=1}^{K} e^{f_\ell(x)}\Bigr).
\end{equation}
与二分类情形一致，准则 \eqref{eq:10_22} 对错误预测的惩罚只随其错误程度\textbf{线性}增长。（Zhu et al., 2005 将指数损失推广到 $K$ 类，见练习 10.5。）

\subsection{回归的稳健损失}

回归中与「指数损失  二项对数似然」类比的是「\textbf{平方误差  绝对损失}」：
\begin{itemize}
    \item 平方误差 $L(y, f) = (y - f)^2$ 的总体解是 $f(x) = \E(Y \mid x)$；
    \item 绝对损失 $L(y, f) = \abs{y - f}$ 的总体解是 $f(x) = \mathrm{median}(Y \mid x)$；对称误差分布下两者相同。
\end{itemize}
但有限样本下，平方误差在拟合过程中对绝对残差 $\abs{y_i - f(x_i)}$ 大的观测施加\textbf{过重}的权重，因而\textbf{不稳健}：对长尾误差分布、尤其对严重误测的 $y$（离群点），性能严重退化；绝对损失等更稳健准则在这些情形表现好得多。稳健统计文献提出了各种对严重离群点强抵抗、对高斯误差又几乎与最小二乘同样有效的回归损失。

\subsubsection{Huber 损失}

其中最具代表性的是 M-回归的 \textbf{Huber 损失}（Huber, 1964）：
\begin{equation}\label{eq:10_23}
    L(y, f(x)) = \begin{cases} [y - f(x)]^2, & \abs{y - f(x)} \le \delta,\\ 2\delta\abs{y - f(x)} - \delta^2, & \text{否则}. \end{cases}
\end{equation}
它在零附近是二次的（像平方误差）、大残差处是线性的（像绝对损失），兼具二者优点。三种回归损失随残差 $y - f(x)$ 的变化对比如下（$\delta = 1$）：

\begin{center}
\begin{tikzpicture}[>=stealth, font=\small]
    % 坐标轴与刻度（先画，避免被 clip 裁剪）
    \draw[->, thick] (-3.6,0) -- (3.7,0) node[below right] {$y - f(x)$（残差）};
    \draw[->, thick] (0,-0.15) -- (0,8.6) node[above left] {损失};
    \foreach \x in {-3,-2,-1,1,2,3} \draw (\x,0.05) -- (\x,-0.05) node[below] {$\x$};
    \foreach \y in {2,4,6,8} \draw (0.05,\y) -- (-0.05,\y) node[left] {$\y$};
    % 曲线（clip：允许平方误差顶部被截断）
    \begin{scope}
        \clip (-3.6,-0.3) rectangle (3.6,8.5);
        % 平方误差 (y-f)^2
        \draw[green!50!black, thick, domain=-3:3, samples=150] plot (\x, {\x*\x});
        % 绝对损失 |y-f|
        \draw[blue, thick, domain=-3:3, samples=150] plot (\x, {abs(\x)});
        % Huber（δ=1）：|r|<=1 为 r^2，否则 2|r|-1（分段绘制保证在 ±1 处精确衔接）
        \draw[red, very thick, domain=-3:-1, samples=50] plot (\x, {-2*\x - 1});
        \draw[red, very thick, domain=-1:1, samples=100] plot (\x, {\x*\x});
        \draw[red, very thick, domain=1:3, samples=50] plot (\x, {2*\x - 1});
    \end{scope}
    % δ 衔接点标注
    \fill (1,1) circle (1.6pt);
    \fill (-1,1) circle (1.6pt);
    \draw[dashed, gray] (1,0) -- (1,1);
    \draw[dashed, gray] (-1,0) -- (-1,1);
    \node[below, gray] at (1,0) {$\delta$};
    \node[below, gray] at (-1,0) {$-\delta$};
    % 图例
    \draw[green!50!black, thick] (1.65,8.05) -- (2.05,8.05) node[right, green!50!black] {平方误差 $(y-f)^2$};
    \draw[blue, thick] (1.65,7.55) -- (2.05,7.55) node[right, blue] {绝对损失 $|y-f|$};
    \draw[red, very thick] (1.65,7.05) -- (2.05,7.05) node[right, red] {Huber（$\delta = 1$）};
\end{tikzpicture}
\\[2pt]
\small\textbf{图：}三种回归损失函数（对应原书图 10.5）。Huber 损失在 $\abs{y-f} \le \delta$ 内为二次（与平方误差重合）、在 $\abs{y-f} > \delta$ 处变为线性（与绝对损失平行），并在 $\pm\delta$ 处光滑衔接。
\end{center}

\textbf{Huber 的两段在 $\pm\delta$ 处光滑衔接：} 当 $\abs{y - f(x)} = \delta$ 时两段取值都为 $\delta^2$，且内段导数 $2\delta$ 恰等于外段斜率 $2\delta$，故函数与导数都连续。其导数为
\[
    \pdv{L}{f} = \begin{cases} -2\,[y - f(x)], & \abs{y - f(x)} \le \delta,\\ -2\delta \cdot \sgn[y - f(x)], & \text{否则}, \end{cases}
\]
即残差的\textbf{影响力}在 $\abs{y-f} \le \delta$ 内随残差线性增长（像最小二乘）、一旦超过 $\delta$ 便\textbf{饱和为常数} $2\delta$（像绝对损失）——这正是其稳健性的来源：大离群点的梯度贡献被「封顶」，不再随残差二次放大。$\delta$ 是「多少残差视为正常」的阈值，常取 $1.345$（使高斯误差下相对渐近效率约 $95\%$）。

\begin{remark}[稳健性与算法优雅性的权衡]
当稳健性是关切点（数据挖掘尤甚，见 §10.7），回归的平方误差、分类的指数损失从统计角度看不是最优准则。但它们在\textbf{前向分阶段加性建模}中导致优雅的模块化 boosting 算法：平方误差下每步只需把基学习器拟合到当前残差 $y_i - f_{m-1}(x_i)$；指数损失下对输出 $y_i$ 做加权拟合、权重 $w_i = \exp(-y_i f_{m-1}(x_i))$。直接换用其他稳健准则得不到这样简单的可行 boosting 算法；但 §10.10.2 将展示如何基于\textbf{任意可微损失}导出简单优雅的 boosting 算法，从而为数据挖掘构造高度稳健的 boosting 程序。
\end{remark}

\paragraph{既然 AdaBoost 是指数损失，那我能设计一个基于 Huber 的 boost 吗？}

\textbf{能}，而且这正是从「AdaBoost = 指数损失」到「梯度提升 GBM（§10.10）」的演进动因之一。

\textbf{框架层面：任意损失都可以 boost。} Boosting 的本质不是「指数损失」，而是\textbf{前向分阶段加性建模}（Algorithm 10.2）——对\textbf{任意}损失 $L$ 都可定义。指数损失的特殊之处只在：二分类 + 离散弱分类器时该子问题有\textbf{闭式解}（加权误差率 → $\beta_m = \frac12\log\frac{1-\mathrm{err}_m}{\mathrm{err}_m}$）。这是指数损失的「运气」，不是 boosting 的必要条件。

\textbf{经典前向分阶段 + Huber 的困难：} 如上文所述，用 Huber 替换平方误差确实能稳健化提升树，但子问题不再好解——平方误差化简为「把树拟合到残差」；Huber 化简不出残差拟合：给区域 $R_{jm}$ 后，式 \eqref{eq:10_30} 的解是 \textbf{Huber location 估计}（M-估计），需迭代（IRLS）；树诱导（找区域）需用近似准则 \eqref{eq:10_27}。

\textbf{现代解法：梯度提升（§10.10）。} 对任意可微损失，梯度提升用\textbf{负梯度作「伪残差」}，绕过子问题不可解：
\[
    r_{im} = -\Bigl[\frac{\partial L(y_i, f(x_i))}{\partial f(x_i)}\Bigr]_{f = f_{m-1}}.
\]
对 Huber 损失，伪残差为
\[
    r_{im} = \begin{cases} 2\bigl(y_i - f_{m-1}(x_i)\bigr), & \abs{\text{残差}} \le \delta,\\ 2\delta\,\sgn\bigl(y_i - f_{m-1}(x_i)\bigr), & \abs{\text{残差}} > \delta, \end{cases}
\]
即小残差按普通残差处理、大残差被「封顶」缩水到 $\pm 2\delta$（winsorized）——这正是 Huber 提升的稳健机制。每步把树拟合到该伪残差，再在各叶节点线性搜索最优 $\gamma$。现代实现早已内置：LightGBM 的 `objective='huber'`、XGBoost 的 pseudo-Huber 都直接支持。

\textbf{若想做「分类的 Huber boost」：} Huber 原生是\textbf{回归}损失（定义在 $y \in \R$）。用于二分类需写成\textbf{边际损失} $\rho(yf(x))$——即「Huber 化的平方 hinge」：$yf > 1$ 时为零、$yf < -1$ 时线性、中间二次（§10.6 提到的图 12.4，Rosset et al., 2004b）。它与平方损失有相同的总体极小化，但兼具二项偏差、二次损失与 SVM hinge 的优点，同样可用梯度提升实现。

\textbf{一句话总结：} 可以设计，而且已经存在。指数损失让 AdaBoost 闭式、优雅但不稳健；Huber（或二项偏差、绝对损失）让 boosting 稳健，代价是每步子问题变难——连接「损失 → 算法」的桥就是梯度提升的伪残差技巧。

% ============ §10.7 数据挖掘的「现成」程序 ============
\section{数据挖掘的「现成」程序}

预测学习是数据挖掘的重要方面。各种方法各有擅长的情形，但实践中\textbf{事先很难知道}哪个程序对给定问题表现最好。表 10.1 汇总了若干学习方法的特性。

\begin{table}[H]
\centering
\caption{不同学习方法的若干特性（表 10.1）。记号：好 = 表现好，一般 = 表现一般，差 = 表现差。}
\begin{tabular}{lccccc}
\toprule
\textbf{特性} & 神经网络 & SVM 核 & 树 & MARS & k-NN \\
\midrule
自然处理「混合」类型数据 & 差 & 差 & 好 & 好 & 差 \\
处理缺失值 & 差 & 差 & 好 & 好 & 好 \\
对输入空间离群点稳健 & 差 & 差 & 好 & 差 & 好 \\
对输入单调变换不敏感 & 差 & 差 & 好 & 差 & 差 \\
计算可扩展性（大 $N$） & 差 & 差 & 好 & 好 & 差 \\
处理无关输入的能力 & 差 & 差 & 好 & 好 & 差 \\
提取特征线性组合 & 好 & 好 & 差 & 差 & 一般 \\
可解释性 & 差 & 差 & 一般 & 好 & 差 \\
预测能力 & 好 & 好 & 差 & 一般 & 好 \\
\bottomrule
\end{tabular}
\end{table}

工业与商业数据挖掘应用对学习程序的要求尤其苛刻：数据集通常\textbf{非常大}（观测数与变量数都多），计算考量重要；数据通常\textbf{杂乱}——输入是定量、二值、分类变量的混合（后者常有多水平）、缺失值普遍、数值变量的分布常长尾且高度偏斜、常含大量严重误测（离群点）、各预测变量尺度差异巨大。此外，通常只有少量预测变量真正与预测相关，且缺乏可靠领域知识来构造相关特征或过滤无关特征——无关特征的加入会严重损害许多方法的性能。最后，数据挖掘一般需要\textbf{可解释}模型：不仅要给出预测，还要提供对输入联合取值与预测响应之间关系的定性理解。因此，神经网络这类「黑箱」方法在纯粹预测场景（如模式识别）很有用，但对数据挖掘用处小得多。

\textbf{「现成（off-the-shelf）」方法}：可直接应用于数据、无需大量耗时的数据预处理或对学习程序的仔细调参。在所有著名学习方法中，\textbf{决策树最接近}满足数据挖掘「现成」程序的要求：构建相对快、产生可解释模型（若树小）；天然容纳数值与分类预测变量的混合及缺失值；对单个预测变量的（严格单调）变换不变，因而尺度化/变换不成问题、对预测变量离群点免疫；把\textbf{内部特征选择}作为过程的一部分，从而对大量无关预测变量的加入有抵抗力（甚至免疫）。这些性质使决策树成为数据挖掘最流行的方法。

但树有一缺陷妨碍它成为理想工具：\textbf{不准确}——其预测精度很少能与数据能达到的最佳水平相比。如 §10.1 所见，boosting 树能显著（常是戏剧性地）提升其精度，同时保留其大部分数据挖掘所需性质；被牺牲的优点是\textbf{速度、可解释性}，以及（对 AdaBoost 而言）对重叠类分布、尤其对训练数据\textbf{错误标注}的稳健性。\textbf{梯度提升模型（GBM）}是树提升的推广，试图缓解这些问题，成为数据挖掘中准确有效的现成程序。


% ============ §10.9 Boosting 树 ============
\section{Boosting 树}

回归与分类树详见 §9.2。树把所有预测变量联合取值空间划分为\textbf{不相交区域} $R_j,\; j = 1, 2, \ldots, J$（树的终端节点），每个区域赋予常数 $\gamma_j$，预测规则为「$x \in R_j \Rightarrow f(x) = \gamma_j$」。树可形式化表示为
\begin{equation}\label{eq:10_25}
    T(x; \Theta) = \sum_{j=1}^{J} \gamma_j I(x \in R_j),
\end{equation}
参数 $\Theta = \{R_j, \gamma_j\}_1^J$，$J$ 通常视为元参数。参数通过最小化经验风险求得：
\begin{equation}\label{eq:10_26}
    \hat\Theta = \argmin_{\Theta} \sum_{j=1}^{J} \sum_{x_i \in R_j} L(y_i, \gamma_j).
\end{equation}
这是困难的组合优化问题，通常满足于近似次优解。把优化拆成两部分有用：
\begin{itemize}
    \item \textbf{给定 $R_j$ 求 $\gamma_j$}：通常平凡，常取 $\hat\gamma_j = \bar y_j$（落入 $R_j$ 的 $y_i$ 均值）；对误分类损失，$\hat\gamma_j$ 取落入 $R_j$ 观测的\textbf{众数类}。
    \item \textbf{求 $R_j$}：困难部分，用近似解；且求 $R_j$ 也暗含估计 $\gamma_j$。典型策略是\textbf{贪心、自上而下的递归划分}算法。有时需用更光滑、更方便的准则 \eqref{eq:10_26} 的近似来优化 $R_j$：
    \begin{equation}\label{eq:10_27}
        \tilde\Theta = \argmin_{\Theta} \sum_{i=1}^{N} \tilde L\bigl(y_i, T(x_i, \Theta)\bigr),
    \end{equation}
    得到 $\hat R_j = \tilde R_j$ 后，再用原始准则更精确地估计 $\gamma_j$。
\end{itemize}

\paragraph{误分类损失下 $\hat\gamma_j$ 取「众数类」是什么？}

\textbf{众数类（modal class）}就是区域内\textbf{出现次数最多的类别}（统计中的\textbf{众数} mode）。原因是误分类损失 $L(y, \gamma) = I(y \neq \gamma)$（0--1 损失）在区域 $R_j$ 内的经验风险为
\[
    \sum_{x_i \in R_j} I(y_i \neq \gamma) = \underbrace{N_j}_{\text{落入 }R_j\text{ 的观测数}} - \underbrace{\#\{y_i = \gamma\}}_{\text{类别恰为 }\gamma\text{ 的观测数}},
\]
其中 $N_j$ 与 $\gamma$ 无关。因此最小化该风险 $\iff$ \textbf{最大化「类别等于 $\gamma$」的计数} $\iff$ 取出现频率最高的类。

\textbf{直觉}：误分类损失每错一个样本罚 $1$。把一个区域 $R_j$ 的预测定成某个类 $\gamma$ 时，付出的代价 = 区域内「不是 $\gamma$」的样本数。显然应把 $\gamma$ 定成区域内样本最多的那个类，使「被罚错」的样本数最少。若区域内各类计数相同（如各 $50\%$），则任取其一即可。

\textbf{与各损失的最优「位置估计」对应：} 给定区域求最优常数，本质是一个「位置」估计问题，损失函数决定了答案：
\begin{center}
\begin{tabular}{lll}
\toprule
\textbf{损失} & \textbf{最优 $\hat\gamma_j$} & \textbf{统计量} \\
\midrule
平方误差 $[y-\gamma]^2$ & $\bar y_j$（均值） & 期望 \\
绝对损失 $\abs{y-\gamma}$ & $\mathrm{median}$（中位数） & 中位数 \\
误分类损失 $I(y\neq\gamma)$ & \textbf{众数类}（modal class） & 众数 \\
\bottomrule
\end{tabular}
\end{center}
三类损失分别对应「均值 / 中位数 / 众数」这三种中心趋势统计量。

§9.2 对分类树描述了这样的策略：生长树（确定 $R_j$）时用 \textbf{Gini 指数}替代误分类损失。

\textbf{提升树模型}是这种树的\textbf{和}：
\begin{equation}\label{eq:10_28}
    f_M(x) = \sum_{m=1}^{M} T(x; \Theta_m),
\end{equation}
以\textbf{前向分阶段}方式（Algorithm 10.2）诱导。前向分阶段每步需对下一棵树求解
\begin{equation}\label{eq:10_29}
    \hat\Theta_m = \argmin_{\Theta_m} \sum_{i=1}^{N} L\bigl(y_i, f_{m-1}(x_i) + T(x_i; \Theta_m)\bigr),
\end{equation}
给定当前模型 $f_{m-1}(x)$ 求区域集与常数 $\Theta_m = \{R_{jm}, \gamma_{jm}\}_1^{J_m}$。

\paragraph{为什么式 \eqref{eq:10_28} 的提升树模型没有显式权重 $\beta_m$？}

权重\textbf{并没有消失}——它被\textbf{吸收}进树的叶节点常数 $\gamma_{jm}$ 里了。回忆 §10.2 的一般加性展开
\[
    f(x) = \sum_{m=1}^{M} \beta_m\, b(x; \gamma_m),
\]
其中 $\beta_m$ 是展开系数（权重）。对树 $T(x;\Theta_m) = \sum_{j=1}^{J} \gamma_{jm} I(x \in R_{jm})$，若写成 $\beta_m T(x;\Theta_m)$，则每叶项变为 $\beta_m \gamma_{jm}$——两个自由标量的乘积仍是自由标量，可重新标度 $\tilde\gamma_{jm} := \beta_m \gamma_{jm}$ 吸收掉 $\beta_m$。\textbf{只要叶节点值 $\gamma_{jm}$ 不受约束（可取任意实数），显式 $\beta_m$ 就是冗余的}，式 \eqref{eq:10_28} 直接令 $\beta_m = 1$ 不失一般性。

\textbf{权重「藏」在哪：} 每棵树的\textbf{整体幅度}体现在其叶节点值 $\gamma_{jm}$ 上——例如平方误差损失下，$\hat\gamma_{jm}$ 是第 $j$ 区域残差的均值（式 \eqref{eq:10_30}），它同时扮演了「拟合残差」与「最优步长」两个角色。前向分阶段每步加的 $T(x;\Theta_m)$，其幅度已由 (10.30) 的最优位置估计确定，无需额外系数。

\textbf{何时必须显式 $\beta_m$：} 只有当下文的\textbf{缩放分类树}（AdaBoost 情形，$\gamma_{jm} \in \{-1, 1\}$）才需要显式 $\beta_m$——因为叶子被\textbf{限制}在 $\pm 1$，幅度无法由 $\gamma_{jm}$ 携带，只能靠外挂的 $\beta_m$。这正是 §10.9 中「二分类 + 指数损失」一项里写 $\beta_m T(x;\Theta_m)$ 的原因；而一般的（不限幅）提升树则无需此系数。

给定区域 $R_{jm}$，求每区域最优常数通常直截了当：
\begin{equation}\label{eq:10_30}
    \hat\gamma_{jm} = \argmin_{\gamma_{jm}} \sum_{x_i \in R_{jm}} L\bigl(y_i, f_{m-1}(x_i) + \gamma_{jm}\bigr).
\end{equation}
求区域是困难的（比单棵树更难）。少数特殊情形下问题简化：
\begin{itemize}
    \item \textbf{平方误差损失}：式 \eqref{eq:10_29} 的解不比单棵树难——只需当前残差 $y_i - f_{m-1}(x_i)$ 预测得最好的回归树，$\hat\gamma_{jm}$ 是各区域残差均值。
    \item \textbf{二分类 + 指数损失}：分阶段方式产生 AdaBoost 提升分类树（Algorithm 10.1）。若树 $T(x;\Theta_m)$ 限制为\textbf{缩放分类树}（$\beta_m T(x;\Theta_m)$，$\gamma_{jm} \in \{-1, 1\}$），则 §10.4 已证式 \eqref{eq:10_29} 的解是使加权误差率 $\sum_{i=1}^{N} w_i^{(m)} I(y_i \neq T(x_i;\Theta_m))$ 最小的树，权重 $w_i^{(m)} = e^{-y_i f_{m-1}(x_i)}$。不加该限制时，指数损失下 \eqref{eq:10_29} 仍简化为新树的加权指数准则：
    \begin{equation}\label{eq:10_31}
        \hat\Theta_m = \argmin_{\Theta_m} \sum_{i=1}^{N} w_i^{(m)} \exp\bigl[-y_i T(x_i; \Theta_m)\bigr].
    \end{equation}
    可用加权指数损失作分裂准则实现贪心递归划分。给定 $R_{jm}$，可证（练习 10.7）式 \eqref{eq:10_30} 的解是各区域的\textbf{加权对数优势}：
    \begin{equation}\label{eq:10_32}
        \hat\gamma_{jm} = \frac{1}{2}\log\frac{\sum_{x_i \in R_{jm}} w_i^{(m)} I(y_i = 1)}{\sum_{x_i \in R_{jm}} w_i^{(m)} I(y_i = -1)}.
    \end{equation}
    这需要专门的树生长算法；实践中更偏好下面用加权最小二乘回归树的近似。
\end{itemize}

用绝对误差或 Huber 损失 \eqref{eq:10_23} 替代平方误差（回归）、用偏差 \eqref{eq:10_22} 替代指数损失（分类），可使提升树\textbf{稳健化}。但不同于其不稳健的对应物，这些稳健准则不产生简单快速的 boosting 算法。对更一般的损失，给定 $R_{jm}$ 后式 \eqref{eq:10_30} 的解通常是直截了当的（一个简单「位置」估计）：绝对损失下是各区域残差的中位数；其他准则有快速迭代算法，通常其更快「单步」近似已足够。问题在\textbf{树诱导}：对这些更一般损失不存在求解 \eqref{eq:10_29} 的简单快速算法，像 \eqref{eq:10_27} 这样的近似变得必不可少。

% ============ §10.10 通过梯度提升的数值优化 ============
\section{通过梯度提升的数值优化}

对\textbf{任意可微损失准则}求解式 \eqref{eq:10_29} 的快速近似算法，可借助与数值优化的类比导出。用 $f(x)$ 预测训练数据上的 $y$ 的总损失为
\begin{equation}\label{eq:10_33}
    L(f) = \sum_{i=1}^{N} L\bigl(y_i, f(x_i)\bigr).
\end{equation}
目标是关于 $f$ 最小化 $L(f)$，此处 $f(x)$ 约束为树的和 \eqref{eq:10_28}。忽略此约束，最小化 \eqref{eq:10_33} 可视为数值优化
\begin{equation}\label{eq:10_34}
    \hat f = \argmin_{f} L(f),
\end{equation}
其中「参数」$f \in \R^N$ 是逼近函数在 $N$ 个数据点处的取值：
\[
    f = \{f(x_1), f(x_2), \ldots, f(x_N)\}^T.
\]
数值优化过程把 \eqref{eq:10_34} 的解表示为分量向量之和
\[
    f_M = \sum_{m=0}^{M} h_m, \qquad h_m \in \R^N,
\]
其中 $f_0 = h_0$ 是初始猜测，每个后继 $f_m$ 基于当前参数向量 $f_{m-1}$（已诱导更新的和）诱导。不同数值优化方法在如何计算每个增量向量 $h_m$（「步长」）上有所不同。

\subsection{最速下降}

\textbf{最速下降（steepest descent）}取 $h_m = -\rho_m g_m$，其中 $\rho_m$ 是标量、$g_m \in \R^N$ 是 $L(f)$ 在 $f = f_{m-1}$ 处求值的梯度。梯度 $g_m$ 的分量为
\begin{equation}\label{eq:10_35}
    g_{im} = \Bigl[\frac{\partial L(y_i, f(x_i))}{\partial f(x_i)}\Bigr]_{f(x_i) = f_{m-1}(x_i)}.
\end{equation}
步长 $\rho_m$ 是如下问题的解：
\begin{equation}\label{eq:10_36}
    \rho_m = \argmin_{\rho} L(f_{m-1} - \rho g_m),
\end{equation}
随后更新 $f_m = f_{m-1} - \rho_m g_m$ 并进入下一次迭代。最速下降可视为\textbf{非常贪心}的策略——$-g_m$ 是 $\R^N$ 中 $L(f)$ 在 $f = f_{m-1}$ 处下降最快的局部方向。

\paragraph{为什么要取「使损失最小」的步长 $\rho_m$？}

方向已定为负梯度 $-g_m$（损失下降最快的局部方向），剩下只需决定\textbf{走多远}。沿这条直线把 $L$ 看成 $\rho$ 的一元函数 $L(f_{m-1} - \rho g_m)$：对（局部）凸的 $L$，它随 $\rho$ 增大\textbf{先降后升}，谷底就是该方向上的最优步长——此时新点 $f_m = f_{m-1} - \rho_m g_m$ 的损失最小，也就是\textbf{沿这个方向损失下降得最多}（与 $f_{m-1}$ 相比的损失下降量最大）：

\begin{center}
\begin{tikzpicture}[>=stealth, font=\small, scale=0.85]
    % 同心椭圆等高线 L = f1^2 + 4 f2^2（中心为极小点 f*）
    \draw[gray!55, thin] (0,0) ellipse (1.41 and 0.71);   % c=2
    \draw[gray!55, thin] (0,0) ellipse (2.47 and 1.23);   % c=6.08（f_m 所在）
    \draw[gray!55, thin] (0,0) ellipse (4.24 and 2.12);   % c=18.01（f_{m-1} 所在）
    \draw[gray!55, thin] (0,0) ellipse (5.48 and 2.74);   % c=30
    % 坐标轴
    \draw[->] (-6.3,0) -- (6.3,0) node[below right] {$f_1$};
    \draw[->] (0,-3.6) -- (0,3.6) node[above left] {$f_2$};
    \fill (0,0) circle (1.8pt) node[below left] {$f^*$};
    % 一维搜索线（过 f_{m-1}，沿 -g_m 方向延长）
    \draw[red!60!black, dashed, thick] (5.25,3.6) -- (0.35,-3.12);
    \node[red!60!black] at (4.55,2.85) {$-g_m$ 方向};
    % 起点 f_{m-1} 与最优步长点 f_m
    \fill (3.5,1.2) circle (1.8pt) node[above right] {$f_{m-1}$};
    \fill (2.317,-0.422) circle (1.8pt) node[below right] {$f_m$};
    % 负梯度步长箭头
    \draw[blue, very thick, ->] (3.5,1.2) -- (2.317,-0.422)
        node[midway, above left, blue] {$-\rho_m g_m$};
\end{tikzpicture}
\\[2pt]
\small\textbf{图（等高线视角）：}最速下降的一步。从 $f_{m-1}$ 沿负梯度方向 $-g_m$（与等高线垂直）做一维线搜索，$f_m$ 取该直线上损失最小的点（直线与等高线\textbf{相切}处）。梯度只给方向，$\rho_m$ 决定走多远。
\end{center}

\begin{center}
\begin{tikzpicture}[>=stealth, font=\small, scale=1.0]
    \draw[->] (-0.12,0) -- (0.52,0) node[below right] {$\rho$};
    \draw[->] (0,-2.5) -- (0,14.5) node[above left] {$L(f_{m-1}-\rho g_m)$};
    % 沿 -g_m 方向的损失曲线（本例 L = f1^2 + 4 f2^2，f_{m-1}=(3.5,1.2)）
    \draw[red!70!black, thick, domain=0:0.28, samples=150]
        plot (\x, {12.25*(1-2*\x)*(1-2*\x) + 5.76*(1-8*\x)*(1-8*\x)});
    % 起点：rho=0 时损失 = L(f_{m-1})
    \fill (0,18.01) circle (1.8pt) node[left] {$L(f_{m-1})$};
    % 谷底：rho_m
    \draw[dashed] (0.169,0) -- (0.169,6.08);
    \fill (0.169,6.08) circle (1.8pt);
    \node[above right] at (0.169,6.08) {$\rho_m$（损失最小）};
    \node[below] at (0.169,-0.1) {$\rho_m$};
\end{tikzpicture}
\\[2pt]
\small\textbf{图（剖面视角）：}沿 $-g_m$ 方向的损失 $L(f_{m-1}-\rho g_m)$ 随 $\rho$ 变化：先降后升，谷底 $\rho_m$ 处损失最小——即相对旧点 $f_{m-1}$ 的损失下降量在该方向上最大。
\end{center}

\subsection{梯度提升}

前向分阶段提升（Algorithm 10.2）同样是\textbf{贪心}策略：给定当前模型 $f_{m-1}$ 及其拟合 $f_{m-1}(x_i)$，每步选出的树是使 \eqref{eq:10_29} 减少最多的解树。因此树预测 $T(x_i;\Theta_m)$ 与负梯度 \eqref{eq:10_35} 的分量\textbf{类似}。主要区别是树的各分量 $t_m = \{T(x_1;\Theta_m), \ldots, T(x_N;\Theta_m)\}^T$ \textbf{不独立}——它们被约束为某棵 $J_m$ 终端节点决策树的预测，而负梯度是\textbf{无约束}的最大下降方向。分阶段方法中 \eqref{eq:10_30} 的解类似于最速下降中的线搜索 \eqref{eq:10_36}；区别是 \eqref{eq:10_30} 对 $t_m$ 中对应每个独立终端区域的分量分别做线搜索。

若唯一目标是极小化训练损失 \eqref{eq:10_33}，最速下降是首选：梯度 \eqref{eq:10_35} 对任何可微损失都\textbf{易于计算}，而 \eqref{eq:10_29} 对 §10.6 讨论的稳健准则很难求解。遗憾的是梯度 \eqref{eq:10_35} 只在训练点 $x_i$ 处有定义，而最终目标是让 $f_M(x)$ \textbf{泛化}到训练集未出现的新数据。

解决这一困境的办法是：在第 $m$ 次迭代诱导一棵预测 $t_m$ 尽量接近负梯度的树 $T(x;\Theta_m)$。用平方误差度量接近程度：
\begin{equation}\label{eq:10_37}
    \tilde\Theta_m = \argmin_{\Theta} \sum_{i=1}^{N} \bigl(-g_{im} - T(x_i;\Theta)\bigr)^2,
\end{equation}
即\textbf{用最小二乘把树 $T$ 拟合到负梯度值} \eqref{eq:10_35}。§10.9 已述最小二乘决策树诱导有快速算法。\eqref{eq:10_37} 解出的区域 $\tilde R_{jm}$ 虽与 \eqref{eq:10_29} 的真正解 $R_{jm}$ 不完全相同，但通常足够相似以达同样目的（前向分阶段与自上而下树诱导本身都是近似过程）。构造树 \eqref{eq:10_37} 后，各区域常数由 \eqref{eq:10_30} 给出。

\paragraph{式 \eqref{eq:10_37} 中「谁和谁」在比较？}

是\textbf{负梯度 $-g_{im}$（理想目标）}与\textbf{树的预测 $T(x_i;\Theta)$（受限近似）}在比较，用平方误差度量二者的差距：
\begin{itemize}
    \item \textbf{$-g_{im}$（负梯度）}：对每个训练点 $x_i$，它给出「$f(x_i)$ 该往哪个方向、降多少，损失 $L$ 才下降最快」——这是\textbf{自由、理想}的逐点目标，但\textbf{只定义在 $N$ 个训练点上}，无法外推到新数据。
    \item \textbf{$T(x_i;\Theta)$（树的预测）}：候选树在 $x_i$ 处的输出——它是定义在整个输入空间上的分段常数函数，\textbf{可以泛化}到任意新点。
\end{itemize}
于是式 \eqref{eq:10_37} 做的事是：\textbf{把 $-g_{im}$ 当作「回归目标 $y$」、把树当作「回归器」}，用最小二乘拟合一棵使二者尽量接近的树——即「让树去模仿/逼近负梯度这个理想方向」。

\textbf{为什么这能解决困境？} 最速下降的更新 $f_m = f_{m-1} - \rho_m g_m$ 依赖负梯度，但 $g_m$ 只定义在训练点；把树拟合成负梯度的近似 $T(x;\Theta_m) \approx -g_m$ 后，树作为\textbf{可泛化函数}可以评估任意 $x$，于是「在函数空间沿负梯度方向下降一步」变成了「把树拟合到负梯度再累加进 $f_M$」。

\textbf{为什么用最小二乘（而非原损失 $L$）度量接近程度？} 真正想解的是 \eqref{eq:10_29}（找使 $L$ 减少最多的树），但对稳健损失很难求解；而\textbf{最小二乘树诱导有快速算法}（§9.2），且对梯度做平方拟合在数值上等价于沿负梯度方向逼近极小点。因此 \eqref{eq:10_37} 是 \eqref{eq:10_29} 的一个\textbf{实用近似}——解出的 $\tilde R_{jm}$ 与真正最优区域 $R_{jm}$ 不完全相同，但通常足够好。

\textbf{常用损失函数的梯度（表 10.2）：}
\begin{center}
\begin{tabular}{lll}
\toprule
\textbf{设定} & \textbf{损失函数} & $-\partial L(y_i, f(x_i))/\partial f(x_i)$ \\
\midrule
回归 & $\frac{1}{2}[y_i - f(x_i)]^2$ & $y_i - f(x_i)$ \\
回归 & $\abs{y_i - f(x_i)}$ & $\sgn[y_i - f(x_i)]$ \\
回归 & Huber & $y_i - f(x_i)$（$\abs{y_i - f(x_i)} \le \delta_m$）\\
     &  & $\delta_m \sgn[y_i - f(x_i)]$（$\abs{y_i - f(x_i)} > \delta_m$）\\
     &  & 其中 $\delta_m$ 是 $\abs{y_i - f(x_i)}$ 的 $\alpha$ 分位数 \\
分类 & Deviance & 第 $k$ 分量：$I(y_i = \mathcal{G}_k) - p_k(x_i)$ \\
\bottomrule
\end{tabular}
\end{center}
平方误差损失下，负梯度即普通残差 $-g_{im} = y_i - f_{m-1}(x_i)$，故 \eqref{eq:10_37} 本身等价于标准最小二乘提升。绝对误差损失下负梯度是残差的符号，每步把树最小二乘拟合到当前残差的符号。Huber M-回归的负梯度是二者的折中（见表）。分类时损失取多项式偏差 \eqref{eq:10_22}，每步构造 $K$ 棵最小二乘树，每棵 $T_{km}$ 拟合各自负梯度向量 $g_{km}$：
\begin{equation}\label{eq:10_38}
    -g_{ikm} = \Bigl[\frac{\partial L(y_i, f_1(x_i), \ldots, f_K(x_i))}{\partial f_k(x_i)}\Bigr]_{f(x_i) = f_{m-1}(x_i)} = I(y_i = \mathcal{G}_k) - p_k(x_i),
\end{equation}
其中 $p_k(x)$ 由 \eqref{eq:10_21} 给出。每步虽建 $K$ 棵树，但它们通过 \eqref{eq:10_21} 相互关联；二分类（$K = 2$）只需一棵树（练习 10.10）。

\subsection{梯度提升的实现}

Algorithm 10.3 给出回归的通用梯度树提升算法；插入不同损失 $L(y, f(x))$ 得到具体算法。第 1 行初始化到最优常数模型（单终端节点树）；第 2(a) 行计算的负梯度分量称为\textbf{广义/伪残差} $r$。

\textbf{Algorithm 10.3（梯度树提升算法）：}
\begin{enumerate}
    \item 初始化 $f_0(x) = \argmin_{\gamma} \sum_{i=1}^{N} L(y_i, \gamma)$。
    \item 对 $m = 1$ 到 $M$：
    \begin{enumerate}
        \item 对 $i = 1, 2, \ldots, N$ 计算 $r_{im} = -\Bigl[\frac{\partial L(y_i, f(x_i))}{\partial f(x_i)}\Bigr]_{f = f_{m-1}}$。
        \item 拟合回归树到目标 $r_{im}$，得到终端区域 $R_{jm},\; j = 1, 2, \ldots, J_m$。
        \item 对 $j = 1, 2, \ldots, J_m$ 计算 $\gamma_{jm} = \argmin_{\gamma} \sum_{x_i \in R_{jm}} L\bigl(y_i, f_{m-1}(x_i) + \gamma\bigr)$。
        \item 更新 $f_m(x) = f_{m-1}(x) + \sum_{j=1}^{J_m} \gamma_{jm} I(x \in R_{jm})$。
    \end{enumerate}
    \item 输出 $\hat f(x) = f_M(x)$。
\end{enumerate}

分类算法类似：每次迭代 $m$ 对每类用 \eqref{eq:10_38} 重复第 2(a)--(d) 行 $K$ 次，最后得到 $K$ 个（耦合的）树展开 $f_{kM}(x),\, k = 1, \ldots, K$，通过 \eqref{eq:10_21} 产生概率或按 \eqref{eq:10_20} 分类（练习 10.9）。两个基本调节参数是迭代次数 $M$ 与各组成树的大小 $J_m$。

原始实现称为 \textbf{MART}（multiple additive regression trees，多元加性回归树），本书第一版如此称呼，本章许多图由 MART 产生。梯度提升在 R 的 \texttt{gbm} 包（Ridgeway, 1999）中实现并免费可用；另一 R 实现是 \texttt{mboost}（Hothorn and Bühlmann, 2006）；商业实现 TreeNet 来自 Salford Systems。

% ============ §10.11 提升树的合适大小 ============
\section{提升树的合适大小}

历史上 boosting 被视为组合模型（这里是树）的技术，树构建算法被当作产生待组合模型的「原始构件」，每棵树的\textbf{最优大小}在构建时按常规方式（§9.2）单独估计：先诱导一棵很大的（过大）树，再用自底向上过程剪枝到估计的最优终端节点数。这一方法\textbf{隐含假设每棵树都是展开 \eqref{eq:10_28} 中的最后一棵}——除最后一棵树外这显然是很差的假设，结果是树通常\textbf{过大}，尤其在早期迭代，显著损害性能并增加计算。

避免此问题的最简单策略是\textbf{限制所有树同大小}：$J_m = J\ \forall m$，每步诱导一棵 $J$ 终端节点回归树。于是 $J$ 成为整个提升过程的\textbf{元参数}，针对手头数据调整以最大化估计性能。

有用 $J$ 值可通过考虑「目标」函数的性质获得：
\begin{equation}\label{eq:10_39}
    \eta = \argmin_{f} \E_{XY} L(Y, f(X)),
\end{equation}
期望取遍 $(X, Y)$ 的总体联合分布；$\eta(x)$ 是对未来数据预测风险最小的函数，即我们要逼近的对象。$\eta(X)$ 的一个相关性质是坐标变量 $X^T = (X_1, \ldots, X_p)$ 之间的\textbf{交互程度}，由它的 \textbf{ANOVA 展开}刻画：
\begin{equation}\label{eq:10_40}
    \eta(X) = \sum_{j} \eta_j(X_j) + \sum_{jk} \eta_{jk}(X_j, X_k) + \sum_{jkl} \eta_{jkl}(X_j, X_k, X_l) + \cdots.
\end{equation}
\eqref{eq:10_40} 第一项是对单个预测变量 $X_j$ 的函数，这些 $\eta_j(X_j)$ 是在所用损失下联合最佳逼近 $\eta(X)$ 的函数，每个称为 $X_j$ 的「\textbf{主效应}」；第二项是加入主效应后对 $\eta(X)$ 拟合最好的两变量函数，称为变量对 $(X_j, X_k)$ 的\textbf{二阶交互}；第三项代表三阶交互，依此类推。实践中遇到的许多问题\textbf{低阶交互占主导}；此时产生强高阶交互的模型（如大树）精度受损。

\textbf{树的大小限制交互阶数}：基于树的近似的交互水平受 $J$ 限制——不可能出现高于 $J-1$ 阶的交互。由于提升模型对树是加性的 \eqref{eq:10_28}，此限制同样适用于它们。取 $J = 2$（单分裂「decision stump」）产生\textbf{只有主效应}的提升模型，不允许交互；$J = 3$ 还允许两变量交互，依此类推。因此 $J$ 的选择应反映 $\eta(x)$ 主导交互的阶数——这通常未知，但多数情况下偏低。图 10.9 用模拟例 \eqref{eq:10_2} 展示交互阶数（$J$ 的选择）的影响：生成函数是加性的（二次单项式之和），故 $J > 2$ 的提升模型招致不必要方差、测试误差更高。图 10.10 比较提升 stump 找到的坐标函数与真函数。

虽然很多应用中 $J = 2$ 不够，但需要 $J > 10$ 的可能性不大。经验表明 boosting 中 $4 \le J \le 8$ 效果好，且对该范围内具体选择相当不敏感。可在验证集上试几个值取风险最低者微调 $J$，但这很少比直接用 $J \simeq 6$ 有显著改进。

% ============ §10.12 正则化 ============
\section{正则化}

除组成树大小 $J$ 外，梯度提升的另一个元参数是提升迭代次数 $M$。每次迭代通常降低训练风险 $L(f_M)$，故 $M$ 足够大时该风险可任意小；但拟合训练数据过好会导致\textbf{过拟合}、损害未来预测风险。故存在极小化未来风险的最优 $M^*$，与应用相关。方便的估计 $M^*$ 的方法是：在验证样本上监控预测风险作为 $M$ 的函数，取使风险最小的 $M$ 作为 $M^*$ 的估计——这与神经网络常用的\textbf{早停}策略类似（§11.4）。

\subsection{收缩}

控制 $M$ 不是唯一的正则化策略；与岭回归和神经网络类似，也可用\textbf{收缩}（shrinkage，§3.4.1、§11.5）。boosting 中最简单的收缩实现：把每棵树的贡献缩放因子 $0 < \nu < 1$ 再加到当前逼近上，即把 Algorithm 10.3 第 2(d) 行替换为
\begin{equation}\label{eq:10_41}
    f_m(x) = f_{m-1}(x) + \nu\cdot\sum_{j=1}^{J} \gamma_{jm} I(x \in R_{jm}).
\end{equation}
参数 $\nu$ 可看作控制 boosting 过程的\textbf{学习率}。更小的 $\nu$（更多收缩）在相同迭代数 $M$ 下导致更大训练风险，故 $\nu$ 与 $M$ 都控制训练数据上的预测风险；但它们\textbf{不独立}——更小 $\nu$ 对应相同训练风险需要更大 $M$，二者存在\textbf{权衡}。

经验上（Friedman, 2001）更小的 $\nu$ 偏好更好的测试误差、需要相应更大的 $M$。事实上最佳策略似乎是设 $\nu$ 很小（$\nu < 0.1$）再用早停选 $M$：对回归与概率估计较无收缩（$\nu = 1$）有显著改进；对 \eqref{eq:10_20} 的误分类风险改进较小但依然可观。代价是计算：更小 $\nu$ 需要更大 $M$，计算量正比于后者；但因每步诱导的树小且不剪枝，即便非常大的数据集上许多迭代通常也可行。图 10.11 用模拟例 \eqref{eq:10_2} 展示测试误差曲线（MART + 二项偏差，stump 或 6 终端节点树，有/无收缩）：收缩的益处明显，尤其跟踪二项偏差时——有收缩时每条测试误差曲线达到更低值且保持多个迭代。§16.2.1 将建立 boosting 前向分阶段收缩与正则化模型参数的 $\ell_1$ 惩罚（lasso）的联系，并论证 $\ell_1$ 惩罚可能优于 SVM 等使用的 $\ell_2$ 惩罚。

\subsection{子采样}

§8.7 已见 bootstrap 平均（bagging）通过平均改进噪声分类器性能；第 15 章详述这种「采样后平均」的降方差机制。梯度提升中可用同一装置\textbf{同时改进性能与计算效率}。\textbf{随机梯度提升}（Friedman, 1999）：每次迭代从训练观测中\textbf{不放回}采样比例 $\eta$，用该子样本生长下一棵树，其余算法相同。$\eta$ 的典型值可为 $\frac12$，$N$ 大时 $\eta$ 可远小于 $\frac12$。采样不仅把计算时间减少同比例 $\eta$，许多情况下还实际产生更准确的模型。图 10.12 用模拟例 \eqref{eq:10_2} 展示子采样效果（分类与回归两例）：两例中「采样 + 收缩」略优于其余；此处「子采样而无收缩」表现不佳。代价是现在要设\textbf{四个参数}：$J$、$M$、$\nu$、$\eta$。通常一些初步探索确定 $J$、$\nu$、$\eta$ 的合适值，留下 $M$ 作主要参数。

% ============ §10.13 解释 ============
\section{解释}

单棵决策树高度可解释——整个模型可用简单的二维图形（二叉树）完全表示、易于可视化。树的线性组合 \eqref{eq:10_28} 丧失了这一重要特性，须以不同方式解释。

\subsection{预测变量的相对重要性}

数据挖掘应用中输入预测变量很少同等相关：往往只有少数对响应有实质影响，绝大多数无关、删掉也无妨。了解每个输入变量在预测响应中的\textbf{相对重要性或贡献}通常很有用。

对单棵决策树 $T$，Breiman et al. (1984) 提出每个预测变量 $X_\ell$ 的相关性度量：
\begin{equation}\label{eq:10_42}
    \mathcal{I}_\ell^2(T) = \sum_{t=1}^{J-1} \hat\imath_t^2\, I(v(t) = \ell),
\end{equation}
对树的 $J-1$ 个内部节点求和。在每个内部节点 $t$，用一个输入变量 $X_{v(t)}$ 把该节点区域划分为两个子区域，每个子区域内用常数拟合响应；被选中的变量是在整个区域用常数拟合之上给出\textbf{最大平方误差风险改进} $\hat\imath_t^2$ 的那个。变量 $X_\ell$ 的平方相对重要性就是「以它作分裂变量的所有内部节点的平方改进之和」。

该度量容易推广到加性树展开 \eqref{eq:10_28}——只需对树求平均：
\begin{equation}\label{eq:10_43}
    \mathcal{I}_\ell^2 = \frac{1}{M}\sum_{m=1}^{M} \mathcal{I}_\ell^2(T_m).
\end{equation}
由于平均的稳定化作用，此度量比单棵树的 \eqref{eq:10_42} 更可靠；且因收缩（§10.12.1），重要变量被与之高度相关的其他变量\textbf{掩盖}的问题大大减轻。注意 \eqref{eq:10_42}、\eqref{eq:10_43} 指\textbf{平方}相关性，实际相关性是各自的平方根。由于这些度量是相对的，习惯上把最大者赋为 100、其余按比例缩放。

对 $K$ 类分类，诱导 $K$ 个独立模型 $f_k(x),\; k = 1, 2, \ldots, K$，每个是树之和
\begin{equation}\label{eq:10_44}
    f_k(x) = \sum_{m=1}^{M} T_{km}(x).
\end{equation}
此时 \eqref{eq:10_43} 推广为
\begin{equation}\label{eq:10_45}
    \mathcal{I}_{\ell k}^2 = \frac{1}{M}\sum_{m=1}^{M} \mathcal{I}_\ell^2(T_{km}),
\end{equation}
其中 $\mathcal{I}_{\ell k}$ 是 $X_\ell$ 在「把第 $k$ 类观测与其他类分开」中的相关性。$X_\ell$ 的\textbf{总体}相关性通过对所有类平均得到：
\begin{equation}\label{eq:10_46}
    \mathcal{I}_\ell^2 = \frac{1}{K}\sum_{k=1}^{K} \mathcal{I}_{\ell k}^2.
\end{equation}

\subsection{部分依赖图}

识别出最相关变量后，下一步是理解逼近 $f(X)$ 对它们\textbf{联合取值}的依赖性质。把 $f(X)$ 作为其自变量的函数图形化，可对输入变量联合取值上的依赖给出全面概括。但可视化限于\textbf{低维}视图——一元或二元函数（连续、离散或混合）易展示；稍高维可通过固定除一两个自变量外的其余值、生成 trellis 图（Becker et al., 1996，R 中 lattice）展示；超过两三个变量则更难。

一个有用的替代：查看一组图，每张显示逼近 $f(X)$ 对选定的\textbf{小变量子集}的部分依赖。虽然这类图通常无法全面刻画逼近，但常能提供有价值的线索，尤其当 $f(x)$ 由低阶交互 \eqref{eq:10_40} 主导时。

考虑输入预测变量 $X^T = (X_1, \ldots, X_p)$ 的子向量 $X_S$（$|\mathcal{S}| = \ell < p$，索引集 $\mathcal{S} \subset \{1, \ldots, p\}$）。设 $\mathcal{C}$ 为补集，$\mathcal{S} \cup \mathcal{C} = \{1, \ldots, p\}$。一般函数 $f(X)$ 原则上依赖所有输入：$f(X) = f(X_S, X_{\mathcal{C}})$。$f(X)$ 对 $X_S$ 的\textbf{平均或部分依赖}定义为
\begin{equation}\label{eq:10_47}
    f_S(X_S) = \E_{X_{\mathcal{C}}} f(X_S, X_{\mathcal{C}}),
\end{equation}
即 $f$ 的\textbf{边际平均}，当 $X_S$ 与 $X_{\mathcal{C}}$ 无强交互时可作为选定子集对 $f(X)$ 影响的有用描述。部分依赖函数可用于解释任何「黑箱」学习方法的输出，可由下式估计：
\begin{equation}\label{eq:10_48}
    \bar f_S(X_S) = \frac{1}{N}\sum_{i=1}^{N} f(X_S, x_{i\mathcal{C}}),
\end{equation}
其中 $\{x_{1\mathcal{C}}, \ldots, x_{N\mathcal{C}}\}$ 是训练数据中出现的 $X_{\mathcal{C}}$ 值。这需要对每个待求 $\bar f_S(X_S)$ 的联合取值遍历一次数据，即使中等规模数据也可能计算密集。幸运的是对决策树，$\bar f_S(X_S)$ \eqref{eq:10_48} 可\textbf{直接从树本身快速计算}、无需引用数据（练习 10.11）。

\begin{remark}[部分依赖 vs 条件期望]
部分依赖 \eqref{eq:10_47} 表示在\textbf{考虑}（平均掉）其他变量 $X_{\mathcal{C}}$ 对 $f$ 的影响之后，$X_S$ 对 $f$ 的效应；它\textbf{不是}忽略 $X_{\mathcal{C}}$ 效应的条件期望
\begin{equation}\label{eq:10_49}
    \tilde f_S(X_S) = \E\bigl(f(X_S, X_{\mathcal{C}}) \mid X_S\bigr),
\end{equation}
后者是仅用 $X_S$ 函数对 $f(X)$ 的最佳最小二乘逼近。$\tilde f_S(X_S)$ 与 $\bar f_S(X_S)$ 仅当 $X_S$ 与 $X_{\mathcal{C}}$ 独立（罕见）时相同。例如若选定子集的效应纯加性
\begin{equation}\label{eq:10_50}
    f(X) = h_1(X_S) + h_2(X_{\mathcal{C}}),
\end{equation}
则 \eqref{eq:10_47} 给出 $h_1(X_S)$（差一个可加常数）；若效应纯乘性
\begin{equation}\label{eq:10_51}
    f(X) = h_1(X_S)\cdot h_2(X_{\mathcal{C}}),
\end{equation}
则 \eqref{eq:10_47} 给出 $h_1(X_S)$（差一个乘法常数）。反之 \eqref{eq:10_49} 两种情形都给不出 $h_1(X_S)$——事实上 \eqref{eq:10_49} 甚至可能对 $f(X)$ \textbf{完全无关}的变量子集产生强效应。
\end{remark}

对提升树逼近 \eqref{eq:10_28} 的选定变量子集查看部分依赖图，有助于对其性质给出定性描述。受计算机图形与人眼感知所限，子集 $X_S$ 的大小必须小（$\ell \approx 1, 2, 3$）。这样的子集数量众多，但只有从通常小得多的「高相关预测变量」集合中选出的那些才可能提供信息；且对 $f$ 效应近似加性 \eqref{eq:10_50} 或乘性 \eqref{eq:10_51} 的子集最具揭示性。

对 $K$ 类分类，有 $K$ 个独立模型 \eqref{eq:10_44}（每类一个），每个通过下式与各自概率 \eqref{eq:10_21} 联系：
\begin{equation}\label{eq:10_52}
    f_k(X) = \log p_k(X) - \frac{1}{K}\sum_{l=1}^{K}\log p_l(X).
\end{equation}
故每个 $f_k(X)$ 是其相应概率在\textbf{对数尺度}上的单调增函数。对每个 $f_k(X)$ \eqref{eq:10_44} 在其最相关预测变量 \eqref{eq:10_45} 上做部分依赖图，有助于揭示实现该类别的对数优势如何依赖相应输入变量。

\subsection{解释工具的谱系：树特有 vs 模型无关}

部分依赖图（\eqref{eq:10_47}/\eqref{eq:10_48}）本身是\textbf{模型无关}的——它对任何「能对新点预测」的黑箱都适用（神经网络、SVM、随机森林、自定义模型皆可）。但它通常和\textbf{重要性度量}配套使用（先定位关键变量、再看其形状），而重要性度量则有树特有与通用之分：

\begin{center}
\begin{tikzpicture}[>=stealth, font=\small,
    root/.style={draw, rounded corners=2pt, inner sep=5pt, align=center, fill=blue!6, thick},
    treebox/.style={draw, rounded corners=2pt, inner sep=4pt, align=center, fill=orange!12},
    genbox/.style={draw, rounded corners=2pt, inner sep=4pt, align=center, fill=green!12},
    shapebox/.style={draw, rounded corners=2pt, inner sep=4pt, align=center, fill=blue!6}]
  % 根：解释工具
  \node[root] (root) at (0,0) {模型解释工具};
  % 左侧：重要性度量（两大类）
  \node[treebox] (tree) at (-3.4,-2.0) {树特有重要性\\分裂改进（\eqref{eq:10_42}/\eqref{eq:10_43}）};
  \node[genbox] (gen) at (2.7,-2.0) {模型无关重要性\\置换重要性、SHAP 等};
  % 右侧：形状刻画
  \node[shapebox] (shape) at (8.0,-2.0) {形状刻画\\部分依赖图（\eqref{eq:10_47}/\eqref{eq:10_48}）\\trellis 图};
  % 箭头
  \draw[->, thick] (root) -- (tree);
  \draw[->, thick] (root) -- (gen);
  \draw[->, thick] (root) -- (shape);
  % 下方说明：配套使用
  \node[align=center, font=\small] at (0,-3.6)
    {定位关键变量（重要性） + 刻画效应形状（部分依赖）\\二者配套，但重要性度量有树特有与通用之分};
\end{tikzpicture}
\\[2pt]
\small\textbf{图：}解释工具的谱系。重要性度量负责「哪些变量重要」——树方法用分裂改进（\eqref{eq:10_42}/\eqref{eq:10_43}，树特有），其他模型用置换重要性、SHAP 等模型无关方法；部分依赖图负责「效应长什么样」，模型无关、任何黑箱都适用。
\end{center}

\begin{itemize}
    \item \textbf{树特有（分裂改进）}：单棵树 $I_\ell^2(T)$（\eqref{eq:10_42}）按「以 $X_\ell$ 作分裂变量的内部节点平方误差改进」累加，多树平均得 \eqref{eq:10_43}——它\textbf{依赖树的分裂结构}，只适用于决策树/随机森林/提升树。
    \item \textbf{模型无关（置换/SHAP）}：\textbf{置换重要性}通过打乱某变量后预测误差上升多少衡量其作用；\textbf{SHAP} 从博弈论 Shapley 值出发分配每个变量的贡献——二者只依赖模型预测，\textbf{任意模型适用}。
    \item \textbf{配套使用}：对树模型，用 \eqref{eq:10_43} 选关键变量 + 部分依赖图看形状；对其他黑箱，用置换重要性/SHAP 选变量 + 部分依赖图看形状——\textbf{部分依赖图始终通用}。
\end{itemize}

% ============ §10.14 示例 ============
\section{示例}

本节在几个较大数据集上、按需要选用不同损失函数来展示梯度提升。

\subsection{California Housing}

该数据集（Pace and Barry, 1997，来自 Carnegie-Mellon StatLib 仓库）由加州 20,460 个社区（1990 年人口普查街区组）的聚合数据组成。响应 $Y$ 是每个社区的中位房价（单位 $10^5$ 美元）。预测变量包括：中位收入 MedInc、住房密度（房屋数 House、平均每房居住人数 AveOccup）、位置（经度 longitude、纬度 latitude）、以及反映房屋属性的若干量（平均房间数 AveRooms、平均卧室数 AveBedrms）。共 8 个数值预测变量。

用 MART 拟合梯度提升模型：$J = 6$ 终端节点、学习率 \eqref{eq:10_41} $\nu = 0.1$、Huber 损失。随机分为训练集（80%）与测试集（20%）。图 10.13 展示平均绝对误差
\begin{equation}\label{eq:10_53}
    \mathrm{AAE} = \E\abs{y - \hat f_M(x)}
\end{equation}
作为训练与测试数据上迭代数 $M$ 的函数：测试误差随 $M$ 增大\textbf{单调下降}，早期较快、随后趋于几乎恒定。因此只要 $M$ 不太小，其具体取值并不关键（许多应用如此）；收缩策略 \eqref{eq:10_41} 倾向消除过拟合问题，对较大数据集尤其如此。

800 次迭代后 $\mathrm{AAE} = 0.31$，可与最优常数预测器 $\mathrm{median}\{y_i\} = 0.89$ 相比。用更熟悉的量：该模型的平方复相关系数 $R^2 = 0.84$。Pace and Barry (1997) 用复杂的空间自回归过程（基于邻近社区中位房价预测，其余预测变量作协变量），经变换实验预测 $\log Y$ 得 $R^2 = 0.85$；以 $\log Y$ 为响应时梯度提升对应值为 $R^2 = 0.86$。

图 10.14 展示 8 个预测变量的相对重要性：中位收入最相关；经度、纬度、平均每房居住人数约有其一半相关性；其余影响力较小。图 10.15 展示最相关非位置变量的单变量部分依赖图。注意这些图\textbf{不完全光滑}——这是基于树模型的后果：决策树产生不连续的分段常数模型 \eqref{eq:10_25}，树之和 \eqref{eq:10_28} 继承之（只是片段更多）。与本书多数方法不同，结果\textbf{不施加光滑约束}，可建模任意尖锐的不连续；曲线总体呈光滑趋势只是因为那是估计出对预测最好的形状。图中底部的刻度线是各变量数据分布的分位数；边缘数据密度较低（尤其较大值处）使曲线在该区域略欠确定。各图纵轴尺度相同，可视觉比较不同变量的相对重要性。

中位房价对中位收入的部分依赖单调递增、主体近似线性；对平均每房居住人数总体单调递减（除小于 1 的情形）；对平均房间数\textbf{非单调}——约 3 个房间处最小、更小与更大都上升。房价对房龄的部分依赖很弱，与其重要性排名 \eqref{eq:10_46}（图 10.14）不一致，提示该弱主效应可能掩盖与其他变量更强的交互。图 10.16 展示房价对房龄与平均每房居住人数联合值的双变量部分依赖：二者明显交互——平均每房居住人数大于 2 时房价几乎不依赖房龄，小于 2 时对房龄强依赖。图 10.17 是拟合模型对经度与纬度联合值的双变量部分依赖（着色等高线图）：中位房价对加州位置有很强的依赖；该图不是「忽略其他预测变量效应的房价对位置图」\eqref{eq:10_49}，而是与其他社区和房屋属性效应\textbf{分离后}位置的影响 \eqref{eq:10_47}，可看作位置溢价——太平洋沿岸（尤其湾区、洛杉矶–圣地亚哥地区）溢价较大，加州北部、中央谷地与东南荒漠地区位置成本低得多。

\subsection{New Zealand Fish}

动植物生态学家用回归模型以环境变量预测物种存在、丰度与丰富度。近年兴趣从简单线性/参数模型转向更复杂模型：GAM（§9.1）、MARS（§9.4）与提升回归树（Leathwick et al., 2005, 2006）。此处建模新西兰周边海域鱼种 Black Oreo Dory 的存在与丰度。

图 10.18 展示 17,000 次拖网（深水网捕鱼，最深 2 km）位置，红点为 Black Oreo 出现的 2,353 次拖网（占经常记录的百余种之一）。每次拖网记录各鱼种捕获量（kg）及若干环境测量：平均拖网深度 AvgDepth、水温与盐度（二者与深度强相关，故 Leathwick et al. (2006) 改用 TempResid 与 SalResid——二者对深度调整后的残差）、海表温度梯度 SSTGrad、反映生态系统生产力的卫星图像指标 Chla、沿海悬浮颗粒物 SusPartMatter（卫星测得）。

分析目标是估计单次拖网发现 Black Oreo 的概率及期望捕获量（对拖网速度、距离、网目尺寸的变化标准化）。作者用 logistic 回归估计概率。对捕获量，自然做法是假设 Poisson 分布并对均值对数建模，但因零过多常不合适；尽管有零膨胀 Poisson（Lambert, 1992）等专门方法，他们选了更简单途径：若 $Y$ 为（非负）捕获量，
\begin{equation}\label{eq:10_54}
    \E(Y \mid X) = \E(Y \mid Y > 0, X)\cdot \Pr(Y > 0 \mid X),
\end{equation}
第二项由 logistic 回归估计，第一项仅用捕获为正的 2,353 次拖网估计。

logistic 回归用梯度提升模型（GBM），二项偏差损失、深度 10 树、收缩因子 $\nu = 0.025$；正捕获回归对 $\log Y$ 用平方误差损失 GBM（同样深度 10 树、$\nu = 0.01$），再对预测取指数还原。两处都用 10 折交叉验证选择项数与收缩因子。

图 10.19（左）展示 GBM 模型序列的平均二项偏差（10 折 CV 与测试数据），比用每项 8 自由度光滑样条拟合的 GAM 略有改进；右图为两模型的 ROC 曲线（§9.2.5）：性能看起来很接近，GBM 或稍有优势（AUC 略高）——等灵敏度/特异度点处 GBM 达 91%、GAM 90%。图 10.20 汇总 logistic GBM 拟合的变量贡献：Black Oreo 有明显深度范围、更常在冷水捕获。所用预测变量均在精细地理网格上可得（源自环境图集、卫星图像），因此可在网格上预测并导入 GIS 制图系统；图 10.21 展示存在性与捕获量的预测图。因其能建模交互、自动选择变量、且对离群点与缺失数据稳健，GBM 在该数据丰富且热忱的生态学界正迅速流行。

\subsection{Demographics Data}

本节用 MART 在多类分类问题上展示梯度提升。数据来自旧金山湾区 9,243 份购物中心顾客问卷（Impact Resources, Inc., Columbus, OH），含 14 个人口统计学问题。此例目标是用其余 13 个变量预测职业、识别区分不同职业类别的人口统计变量。随机分为训练集（80%）与测试集（20%），用 $J = 6$ 节点树、学习率 $\nu = 0.1$。

图 10.22 展示 $K = 9$ 个职业类别及其误差率：总体误差率 $42.5\%$，可比的是「预测最多类别 Prof/Man（专业/管理）」的零模型率 $69\%$。预测最好的四类是 Retired、Student、Prof/Man、Homemaker。图 10.23 展示对所有类平均的预测变量相对重要性 \eqref{eq:10_46}；图 10.24 展示四个最好预测类别各自的相对重要性分布 \eqref{eq:10_45}——最相关预测变量通常随类别而不同，例外是 age（预测 Retired、Student、Prof/Man 时都排前三）。图 10.25 展示这三类 log-odds \eqref{eq:10_52} 对 age 的部分依赖：分离其他变量贡献后，老年人退休对数优势更高、学生相反；专业/管理人员对数优势在中年人最高。逐类检查部分依赖可得到合理结果。

% ============ 本章总结 ============
\section{本章总结}

\textbf{Boosting} 是本章的主题：它把「弱学习器」（通常是小树）顺序地组合成强模型。本章的叙事可以概括为一条「从算法到理论再到工具」的链条。

\begin{center}
\begin{tikzpicture}[>=stealth, font=\small,
    box/.style={draw, rounded corners=2pt, inner sep=5pt, align=center, fill=blue!6},
    theory/.style={draw, rounded corners=2pt, inner sep=5pt, align=center, fill=orange!12},
    tool/.style={draw, rounded corners=2pt, inner sep=5pt, align=center, fill=green!12},
    arr/.style={->, thick}]
  % 第一层：算法
  \node[box] (a1) at (0,0) {AdaBoost.M1\\（Algorithm 10.1）};
  \node[box] (a2) at (4.6,0) {前向分阶段\\加性建模（10.2）};
  \node[box] (a3) at (9.2,0) {梯度提升 GBM\\（Algorithm 10.3）};
  % 第二层：理论
  \node[theory] (t1) at (0,-2.6) {指数损失\\$e^{-yf}$};
  \node[theory] (t2) at (4.6,-2.6) {任意可微损失\\+ 负梯度（伪残差）};
  \node[theory] (t3) at (9.2,-2.6) {稳健损失：二项偏差、Huber};
  % 第三层：工具
  \node[tool] (u1) at (0,-4.0) {加权投票 $G(x)=\sgn(\sum \alpha_m G_m)$};
  \node[tool] (u2) at (4.6,-5.2) {拟合残差/伪残差 + 叶节点常数};
  \node[tool] (u3) at (9.2,-4.0) {正则化：$J,\ M,\ \nu,\ \eta$\\+ 解释：重要性、部分依赖};
  % 垂直箭头：算法 → 理论 → 工具
  \draw[arr] (a1) -- (t1);
  \draw[arr] (a2) -- (t2);
  \draw[arr] (a3) -- (t3);
  \draw[arr] (t1) -- (u1);
  \draw[arr] (t2) -- (u2);
  \draw[arr] (t3) -- (u3);
  % 水平箭头：从左到右的推广
  \draw[arr, dashed] (a1) -- (a2) node[midway, above] {等价};
  \draw[arr, dashed] (a2) -- (a3) node[midway, above] {推广};
  \draw[arr, dashed] (t1) -- (t2);
  \draw[arr, dashed] (t2) -- (t3);
  \draw[arr, dashed] (u1) -- (u2);
  \draw[arr, dashed] (u2) -- (u3);
\end{tikzpicture}
\\[2pt]
\small\textbf{图：}本章主线——AdaBoost.M1 等价于指数损失下的前向分阶段加性建模；把损失推广到任意可微损失并借助负梯度（伪残差）技巧，得到通用的梯度提升（GBM）；再配以稳健损失、正则化与解释工具，构成数据挖掘的「现成」强学习器。
\end{center}

\subsection{核心结论}

\begin{itemize}
    \item \textbf{Boosting = 前向分阶段拟合加性展开}：$f_M(x) = \sum_m \beta_m b(x;\gamma_m)$，顺序添加基学习器（通常是小树）、不回改已添加项。
    \item \textbf{AdaBoost 的等价性}：AdaBoost.M1 恰是\textbf{指数损失} $e^{-yf}$ 下的前向分阶段加性建模；指数损失的总体极小化是类概率对数优势的一半。
    \item \textbf{损失的选择}：指数损失不稳健（对标签噪声/大负边际极敏感）；二项偏差、Huber 等稳健损失让 boosting 更可靠，但失去闭式解——用\textbf{梯度提升}（负梯度作伪残差）处理任意可微损失。
    \item \textbf{梯度提升的实质}：每步把树最小二乘拟合到当前负梯度（伪残差），再做叶节点最优常数；它是「函数空间的最速下降」被约束到树的可泛化近似。
    \item \textbf{正则化三件套}：树的规模 $J$（限制交互阶数，通常 $4\le J\le 8$）、迭代数 $M$（早停）、学习率 $\nu$（收缩）、子采样 $\eta$（随机梯度提升）——四参数协同控制偏差-方差。
    \item \textbf{解释}：相对重要性（\eqref{eq:10_43}）定位关键变量，部分依赖图（\eqref{eq:10_48}）刻画其形状与交互——树方法独有的快速计算使解释高效可行。
    \item \textbf{定位}：boosting 是数据挖掘的「现成」强学习器，擅长混合类型数据、缺失值、无关变量与离群点，常显著优于单棵树与 GAM 等。
\end{itemize}

